% ==================================================================
% 令和9年度入試対策　新潟大学理学部物理学プログラム　基礎学力試験（学校推薦型選抜Ⅰ型）
% 予想練習問題（問題編・解答解説編）
%   セット1・2：各 120 分・300 点（4 大問×75 点）／ セット3（追加）：力学・電気中心のやや難
% デザイン：07「数物tex」原本仕様のプリアンブルを、そのまま使用（変更は次の2点のみ）
%   1) ヘッダー表記「数学解説」→「物理解説」
%   2) 物理の図（tikzify雛形 physics/circuit/plot/geometry）に必要な
%      TikZ ライブラリと circuitikz・pgfplots の追加
% コンパイル：LuaLaTeX（日本語は jlreq / luatexja）。.sty の同梱は不要。
%   必要：tcolorbox(most)・multicolrule・circuitikz・pgfplots・enumitem・fancyhdr ほか
%   ヘッダー罫線は remember picture を使うため、2 回コンパイルすること。
% 状態：作成環境に jlreq/luatexja が無く、このファイルは【未コンパイル】。
%   構造・図・数式は、XeLaTeX の簡易環境（日本語フォントのみ代用）で試験した。
% ==================================================================
\documentclass[paper=b5, fontsize=8pt, fleqn, twoside]{jlreq}
% --- 基本パッケージ ---
\usepackage{luatexja, multicol, amsmath, amssymb, fancyhdr, enumitem, calc, varwidth}
% --- 独自パッケージ（環境にない場合はコメントアウトのままにする） ---
%\usepackage{luwa-ul} 
%\usepackage{KKsymbols}
%\usepackage{keisennote}
% --- TikZ & Graphics ---
\usepackage{tikz}
\usetikzlibrary{shapes.geometric, shapes.arrows, calc} % dart形状などのために必須
% --- 物理の図（tikzify雛形: physics / circuit / plot / geometry）に必要な追加 ---
\usetikzlibrary{arrows.meta, patterns.meta, decorations.pathmorphing, angles, quotes}
\usepackage{circuitikz}
\usepackage{pgfplots}
\pgfplotsset{compat=1.18}
\usepackage[tikz]{multicolrule}
\SetMCRule{line-style=dense-solid-circles, width=0.8pt} 
% --- TColorBox ---
\usepackage[most]{tcolorbox}
% --- フォント・寸法設定 ---
\makeatletter
\def\ascb@textgt#1{\textgt{#1}}
\def\ascb@gtfamily{\gtfamily}
\def\ascb@zw#1#2{#1\zw}
\newdimen\ascb@parindent@dimen
\ascb@parindent@dimen=\zw
\newcommand{\ascb@parindent}[1]{\setlength{\parindent}{#1}\relax}
\setlength{\parindent}{1\zw}
% --- Custom Box: simple (参考問題用) ---
\DeclareTColorBox{simple}{ o m O{.5} O{} }{
  empty, left=2mm, right=2mm, top=-1mm, 
  attach boxed title to top left={xshift=\ascb@zw{1.2}{11pt}}, 
  boxed title style={empty,left=-.5mm,right=-.5mm}, 
  colframe=black, coltitle=black, coltext=black, breakable, 
  before upper={\ascb@parindent{\ascb@parindent@dimen}},
  underlay unbroken={\draw[black,line width=#3pt](title.east) -- (title.east-|frame.east) -- (frame.south east) -- (frame.south west) -- (title.west-|frame.west) -- (title.west); },
  underlay first={\draw[black,line width=#3pt](title.east) -- (title.east-|frame.east) -- (frame.south east) ; \draw[black,line width=#3pt] (frame.south west) -- (title.west-|frame.west) -- (title.west); },
  underlay middle={\draw[black,line width=#3pt](frame.north east) -- (frame.south east) ; \draw[black,line width=#3pt](frame.south west) -- (frame.north west) ;},
  underlay last={\draw[black,line width=#3pt](frame.north east) -- (frame.south east) -- (frame.south west) -- (frame.north west) ;},
  fonttitle=\ascb@gtfamily, IfValueTF={#1}{title=【#2】〈#1〉}{title=【#2】}, #4
}
% --- Custom Box: ptbs (重要事項・KEY) ---
\newlength{\len@ptbs@kk@D}\newlength{\lenn@ptbs@kk@D}\newlength{\myfontsize@ptbs@kk@D}
\newcommand{\titlelength@ptbs@kk@D}[1]{
  \setlength{\myfontsize@ptbs@kk@D}{\f@size pt}
  \def\titletext@ptbs@D{\gtfamily\normalsize\selectfont#1}
  \settowidth{\lenn@ptbs@kk@D}{\titletext@ptbs@D}
  \setlength{\len@ptbs@kk@D}{\linewidth}
  \addtolength{\len@ptbs@kk@D}{-\lenn@ptbs@kk@D}
  \addtolength{\len@ptbs@kk@D}{-.3em}\addtolength{\len@ptbs@kk@D}{-1pt}\addtolength{\len@ptbs@kk@D}{4mm}
}
\DeclareTColorBox{ptbs}{ m O{\phantom{A}} O{} }{
  enhanced, breakable, boxsep=0mm, lefttitle=1.5mm,
  arc=.5mm, bottom=2mm, top=2mm, leftupper=4mm, rightupper=4mm,
  colbacktitle=black!100!white, colframe=black!100!white,
  coltitle=white, colback=black!10!white, boxrule=1pt, lefttitle=.3em,
  before upper={\ascb@parindent{\ascb@parindent@dimen}},
  fonttitle=\gtfamily\normalsize, fontupper=\gtfamily\normalsize,
  title={\titlelength@ptbs@kk@D{#1}#1\kern.3\zw\kern1pt},  
  after title={\tcbox[on line, boxsep=.25\myfontsize@ptbs@kk@D, boxrule=0pt, top=.1\myfontsize@ptbs@kk@D, bottom=.1\myfontsize@ptbs@kk@D, left= .5mm, right=.5mm, width=\len@ptbs@kk@D, colback=black!30!white, arc=.5mm]{\raisebox{.2ex}{\parbox{\len@ptbs@kk@D-.5\myfontsize@ptbs@kk@D-1mm}{\renewcommand{\baselinestretch}{.5}\selectfont#2}}}}, #3
}
% --- Custom Box: ascolorbox4A (問題文用・装飾付き) ---
\DeclareTColorBox{ascolorbox4A}{ o m O{3} O{}}{
  enhanced, colback=white, colframe=white,
  attach boxed title to top left={xshift=1cm,yshift=-\tcboxedtitleheight/2}, 
  varwidth boxed title=0.85\linewidth, coltitle=black, 
  fonttitle=\ascb@gtfamily, before skip=.5mm, after skip=.8mm,
  before upper={\ascb@parindent{\ascb@parindent@dimen}},
  enlarge top by=2mm, enlarge bottom by=2mm, breakable, sharp corners,
  boxed title style={colback=white,left=-.6em,right=-.6em}, 
  borderline={.75pt}{#3pt}{black,dotted},
  underlay unbroken={
    \draw[black,line width=.5pt] (title.east|-frame.north east)--([xshift=-#3*4pt]frame.north east) arc [start angle=180, end angle=270, radius=#3*4pt] -- ([yshift=#3*4pt]frame.south east) arc [start angle=90, end angle=180, radius=#3*4pt] -- ([xshift=#3*4pt]frame.south west) arc [start angle=0, end angle=90, radius=#3*4pt] -- ([yshift=-#3*4pt]frame.north west) arc [start angle=270, end angle=360, radius=#3*4pt] -- (frame.north west-|title.west) ;
    \filldraw[fill=gray,draw=gray] (frame.north east) -- ++(0,-#3*3pt) arc [start angle=270, end angle=180, radius=#3*3pt] -- cycle ;
    \filldraw[fill=gray,draw=gray] (frame.north west) -- ++(#3*3pt,0) arc [start angle=0, end angle=-90, radius=#3*3pt] -- cycle ;
    \filldraw[fill=gray,draw=gray] (frame.south west) -- ++(0,#3*3pt) arc [start angle=90, end angle=0, radius=#3*3pt] -- cycle ;
    \filldraw[fill=gray,draw=gray] (frame.south east) -- ++(-#3*3pt,0) arc [start angle=180, end angle=90, radius=#3*3pt] -- cycle ;
  },
  underlay first={
    \draw[black,line width=.5pt] (title.east|-frame.north east)--([xshift=-#3*4pt]frame.north east) arc [start angle=180, end angle=270, radius=#3*4pt] -- (frame.south east) ;
    \draw[black,line width=.5pt] (frame.south west) -- ([yshift=-#3*4pt]frame.north west) arc [start angle=270, end angle=360, radius=#3*4pt] -- (frame.north west-|title.west) ;
    \filldraw[fill=gray,draw=gray] (frame.north east) -- ++(0,-#3*3pt) arc [start angle=270, end angle=180, radius=#3*3pt] -- cycle ;
    \filldraw[fill=gray,draw=gray] (frame.north west) -- ++(#3*3pt,0) arc [start angle=0, end angle=-90, radius=#3*3pt] -- cycle ;
  },
  underlay middle={
    \draw[black,line width=.5pt] (frame.north east)--(frame.south east) ;
    \draw[black,line width=.5pt] (frame.south west)--(frame.north west) ;
  },
  underlay last={
    \draw[black,line width=.5pt] (frame.north east) -- ([yshift=#3*4pt]frame.south east) arc [start angle=90, end angle=180, radius=#3*4pt] -- ([xshift=#3*4pt]frame.south west) arc [start angle=0, end angle=90, radius=#3*4pt] -- (frame.north west) ;
    \filldraw[fill=gray,draw=gray] (frame.south west) -- ++(0,#3*3pt) arc [start angle=90, end angle=0, radius=#3*3pt] -- cycle ;
    \filldraw[fill=gray,draw=gray] (frame.south east) -- ++(-#3*3pt,0) arc [start angle=180, end angle=90, radius=#3*3pt] -- cycle ;
  },
  IfValueTF={#1}{title=【#2】〈#1〉}{title=【#2】},#4
}
% --- Custom Command: ascboxZ (小見出し用) ---
\tcbset{ascbox@ascolorbox/.style={after skip=1.5mm, before skip=3mm},
ascboxsizeset@ascolorbox/.style={top=0mm,bottom=0mm,right=-1mm,left=2mm,},
titleunderline@ascolorbox/.style={underlay pre={\draw[very thick,draw=gray] ([yshift=.7mm,xshift=3mm]frame.south west) -- ([yshift=.7mm]frame.south east);}}}
\DeclareTCBox{\ascboxZ}{O{dart} s O{.6} s }{
  empty,ascbox@ascolorbox,ascboxsizeset@ascolorbox,
  IfBooleanTF={#4}{}{titleunderline@ascolorbox},
  IfBooleanTF={#2}{underlay={\node[#1,thick,draw=black!40!white,fill=black!70!white,draw,inner sep=#3mm] at (frame.west) {};}}{underlay={\node[#1,thick,draw=black!70!white,fill=black!40!white,draw,inner sep=#3mm] at (frame.west) {};}}
}
% --- Header/Footer (Geometry & TikZ) ---
\newdimen\top@geom@TETSUMANE \top@geom@TETSUMANE=20mm
\newdimen\bottom@geom@TETSUMANE \bottom@geom@TETSUMANE=20mm
\newdimen\left@geom@TETSUMANE \left@geom@TETSUMANE=16mm
\newdimen\right@geom@TETSUMANE \right@geom@TETSUMANE=16mm
\newcommand{\Rhead@TETSUMANE}[1]{\begin{tikzpicture}[remember picture, overlay]
  \draw[line width=.5pt] ([yshift=-\top@geom@TETSUMANE+3mm, xshift=-\left@geom@TETSUMANE] current page.north east) --  ([yshift=-\top@geom@TETSUMANE+3mm, xshift=\right@geom@TETSUMANE] current page.north west);
  \node[anchor=east,yshift=-\top@geom@TETSUMANE+6mm, xshift=-\left@geom@TETSUMANE] at (current page.north east) {\footnotesize #1};
  \draw[line width=.5pt] ([yshift=\bottom@geom@TETSUMANE-3mm, xshift=-\left@geom@TETSUMANE] current page.south east) --  ([yshift=\bottom@geom@TETSUMANE-3mm, xshift=\right@geom@TETSUMANE] current page.south west);
  \node[anchor=center] at ([yshift=\bottom@geom@TETSUMANE-6mm, xshift=0.5*(-\left@geom@TETSUMANE+\right@geom@TETSUMANE)]current page.south) {\footnotesize\gtfamily\symbol{"2015}\hspace*{1ex}\thepage\hspace*{1ex}\symbol{"2015}};
\end{tikzpicture}}
\newcommand{\Lhead@TETSUMANE}[1]{\begin{tikzpicture}[remember picture, overlay]
  \draw[line width=.5pt] ([yshift=-\top@geom@TETSUMANE+3mm, xshift=\left@geom@TETSUMANE] current page.north west) --  ([yshift=-\top@geom@TETSUMANE+3mm, xshift=-\right@geom@TETSUMANE] current page.north east);
  \node[anchor=west,yshift=-\top@geom@TETSUMANE+6mm, xshift=\left@geom@TETSUMANE] at (current page.north west){\footnotesize #1};
  \draw[line width=.5pt] ([yshift=\bottom@geom@TETSUMANE-3mm, xshift=-\right@geom@TETSUMANE] current page.south east) --  ([yshift=\bottom@geom@TETSUMANE-3mm, xshift=\left@geom@TETSUMANE] current page.south west);
  \node[anchor=center] at ([yshift=\bottom@geom@TETSUMANE-6mm,xshift=0.5*(\left@geom@TETSUMANE-\right@geom@TETSUMANE)]current page.south) {\footnotesize\gtfamily\symbol{"2015}\hspace*{1ex}\thepage\hspace*{1ex}\symbol{"2015}};
\end{tikzpicture}}
\usepackage[top=\top@geom@TETSUMANE,bottom=\bottom@geom@TETSUMANE,left=\left@geom@TETSUMANE,right=\right@geom@TETSUMANE]{geometry}
\pagestyle{fancy} \renewcommand{\headrule}{} \fancyhead{} \fancyfoot{}
\fancyhead[ER]{\Rhead@TETSUMANE{\gtfamily 物理解説}}
\fancyhead[OL]{\Lhead@TETSUMANE{\gtfamily 物理解説}}
\makeatother
% --- 余白制御: ページ下端揃えを無効化（必須・削除禁止）---
\raggedbottom
\newcommand{\notefill}{\vfill\null}
\begin{document}
% ------------------------------------------------------------------
% 本文用の補助マクロ（プリアンブル本体は変更せず、ここで定義する）
% ------------------------------------------------------------------
\newcommand{\unit}[1]{\,\mathrm{#1}}                              % 単位は立体・数値との間に薄いスペース
\newcommand{\ans}{\hspace{\zw}\textgt{……(答)}}                      % 答えの行末（文中）
\newcommand{\ansm}{\hspace{\zw}\text{\textgt{……(答)}}}              % 答えの行末（数式中）
\newcommand{\pts}[1]{\nobreak\hfill\textgt{\mbox{［#1点］}}}               % 小問の配点
\newlist{qlist}{enumerate}{1}
\setlist[qlist]{label=\textbf{(\arabic*)},leftmargin=2.4\zw,itemsep=1pt,topsep=2pt,parsep=0pt,partopsep=0pt}
\newcommand{\bigpart}[1]{\par\noindent\textbf{【#1】}\ }
\newcommand{\bpart}[1]{\textbf{【#1】}\ }
% ==================================================================
%  表紙の注記
% ==================================================================
\begin{ptbs}{予想問題}[令和9年度入試対策（新潟大学理学部物理学プログラム）]
学校推薦型選抜Ⅰ型の基礎学力試験を想定した予想練習問題（問題編・解答解説編）。出題範囲は物理基礎（全項目）と物理（力と運動、気体分子の運動、波、電気と磁気）で、「原子」は範囲外のため出題しない。水準は高等学校教科書の章末問題と同じ水準を目標とした。
\par
セット1・2 は各 120 分・300 点（4 大問×75 点）。セット3 は、力学と電気を中心にした、やや難しめの追加セットである。なお、公開過去問（令和6〜8年度）は PDF から本文を読み取れず分析できなかったため、大問数・配点・分野配分は暫定の想定である。
\end{ptbs}
% ==================================================================
%  問題編　セット1
% ==================================================================
\begin{simple}[120分・300点]{予想問題　セット1　注意事項}
\begin{itemize}[leftmargin=1.2\zw,itemsep=0pt,topsep=1pt,parsep=0pt]
\item 4 大問（各 75 点・目標 30 分）。大問は易しいものから難しいものへ並べてある。電卓・定規・コンパスは使えないものとし、図やグラフは徒手で描く。
\item 重力加速度の大きさは $g=9.8\unit{m/s^2}$ とする（文字 $g$ で答える設問は文字のままでよい）。
\item 数値で答える設問は、指定がなければ有効数字 2 桁とし、単位を必ず付ける（割り切れる場合は、そのままの値でよい）。$\sqrt{\phantom{x}}$、$\pi$、分数は計算せずに残してよい。
\item 空気の抵抗は無視する。答えだけでなく、考え方と途中の式も記すこと。
\end{itemize}
\end{simple}

% ---------------------------- 大問1 ----------------------------
\begin{ascolorbox4A}[75点・30分]{大問1　熱と気体}
\bigpart{Ⅰ}熱容量の無視できる断熱容器に、$20\,^\circ\mathrm{C}$ の水 $140\unit{g}$ が入っている。ここに $100\,^\circ\mathrm{C}$ に熱した金属球 $100\unit{g}$ を静かに入れ、十分に時間がたつと全体の温度は $30\,^\circ\mathrm{C}$ になった。水の比熱を $4.2\unit{J/(g\cdot K)}$ とし、熱は水と金属球の間でのみ移動するものとする。
\begin{qlist}[series=sa]
\item 水が得た熱量 $Q$ は何 $\mathrm{J}$ か。\pts{6}
\item 金属の比熱 $c$ は何 $\mathrm{J/(g\cdot K)}$ か。\pts{6}
\item この金属球の熱容量は何 $\mathrm{J/K}$ か。\pts{4}
\end{qlist}
\par\noindent
\begin{minipage}[t]{0.60\linewidth}\vspace{0pt}\bpart{Ⅱ}$n\unit{mol}$ の単原子分子理想気体を、図1のように A→B→C→D→A とゆっくり変化させた。A→B と C→D は定積変化、B→C と D→A は定圧変化である。気体定数を $R$、状態 A の圧力と体積をそれぞれ $p_0$、$V_0$ とする。単原子分子理想気体の内部エネルギーは、絶対温度 $T$ を用いて $U=\dfrac{3}{2}nRT$ と表される。
\end{minipage}\hfill
\begin{minipage}[t]{0.37\linewidth}\vspace{0pt}\centering
\begin{tikzpicture}[>={Stealth[length=1.8mm]},line width=0.6pt,x=1.15cm,y=1.0cm]
  \draw[->] (0,0) -- (3.0,0) node[right] {$V$};
  \draw[->] (0,0) -- (0,2.7) node[above] {$p$};
  \node[below left] at (0,0) {$\mathrm{O}$};
  \draw[densely dashed] (1,0) node[below] {$V_0$} -- (1,1);
  \draw[densely dashed] (2,0) node[below] {$2V_0$} -- (2,1);
  \draw[densely dashed] (0,1) node[left] {$p_0$} -- (1,1);
  \draw[densely dashed] (0,2) node[left] {$2p_0$} -- (1,2);
  \draw[thick] (1,1) -- (1,2) -- (2,2) -- (2,1) -- cycle;
  \draw[thick,->] (1,1.40) -- (1,1.60);
  \draw[thick,->] (1.40,2) -- (1.60,2);
  \draw[thick,->] (2,1.60) -- (2,1.40);
  \draw[thick,->] (1.60,1) -- (1.40,1);
  \fill (1,1) circle[radius=0.05] node[below left] {A};
  \fill (1,2) circle[radius=0.05] node[above left] {B};
  \fill (2,2) circle[radius=0.05] node[above right] {C};
  \fill (2,1) circle[radius=0.05] node[below right] {D};
\end{tikzpicture}\\[1pt]{\small 図1}
\end{minipage}
\begin{qlist}[resume*=sa]
\item 状態 A の絶対温度 $T_\mathrm{A}$ を $n,R,p_0,V_0$ で表せ。また、状態 B, C, D の絶対温度 $T_\mathrm{B},T_\mathrm{C},T_\mathrm{D}$ を $T_\mathrm{A}$ で表せ。\pts{10}
\item 1 サイクルの間に、気体が外部へする正味の仕事 $W$ を $p_0,V_0$ で表せ。\pts{8}
\item 過程 A→B で気体が吸収する熱量 $Q_\mathrm{AB}$ と、過程 B→C で気体が吸収する熱量 $Q_\mathrm{BC}$ を、それぞれ $p_0,V_0$ で表せ。\pts{12}
\item 過程 D→A で、気体は熱を吸収するか、放出するか。仕事と内部エネルギーの変化に着目して、理由とともに答えよ。\pts{10}
\item このサイクルの熱効率 $e$ を、分数で答えよ。\pts{10}
\item 状態 C の気体分子の二乗平均速度は、状態 A の何倍か。\pts{9}
\end{qlist}
\end{ascolorbox4A}

\clearpage
% ---------------------------- 大問2 ----------------------------
\begin{ascolorbox4A}[75点・30分]{大問2　波}
\bigpart{Ⅰ}$x$ 軸の正の向きに進む正弦波がある。図2は時刻 $t=0$ における波形で、縦軸は媒質の変位 $y$ である。波の速さは $v=2.0\unit{m/s}$ である。
\begin{center}
\begin{tikzpicture}
\begin{axis}[width=9.2cm,height=3.7cm,axis lines=middle,
  xlabel={$x\,[\mathrm{m}]$},ylabel={$y\,[\mathrm{cm}]$},
  xlabel style={at={(ticklabel* cs:1.0)},anchor=west},
  ylabel style={at={(ticklabel* cs:1.0)},anchor=south},
  xmin=-0.3,xmax=4.7,ymin=-2.9,ymax=2.9,
  xtick={1,2,3,4},ytick={-2,2},
  tick label style={font=\small},label style={font=\small}]
  \addplot[thick,domain=0:4,samples=100] {2*sin(deg(pi*x/2))};
  \draw[->,thick] (axis cs:0.25,2.55) -- (axis cs:1.15,2.55) node[right] {$v$};
\end{axis}
\end{tikzpicture}\\[1pt]{\small 図2}
\end{center}
\begin{qlist}[series=sb]
\item この波の振幅 $A\,[\mathrm{cm}]$、波長 $\lambda\,[\mathrm{m}]$、周期 $T\,[\mathrm{s}]$、振動数 $f\,[\mathrm{Hz}]$ を求めよ。\pts{8}
\item $t=0$ において、$x=2.0\unit{m}$ にある媒質の速度の向きは、$y$ 軸の正の向きか負の向きか。理由とともに答えよ。\pts{6}
\item 時刻 $t=0.50\unit{s}$ の波形を、図2に重ねて描け。\pts{8}
\item $x=1.0\unit{m}$ にある媒質の変位が、$t=0$ のあとはじめて $0$ となるのは、何秒後か。\pts{6}
\end{qlist}
\par\noindent
\begin{minipage}[t]{0.52\linewidth}\vspace{0pt}\bpart{Ⅱ}図3のように、一端を可動の板で閉じた管の開口端のそばにスピーカーを置き、振動数 $680\unit{Hz}$ の音を出す。板を動かして管の長さ $L$ を $0$ から次第に長くすると、$L=L_1=0.110\unit{m}$ で 1 回目の共鳴が、$L=L_2=0.360\unit{m}$ で 2 回目の共鳴が起こった。開口端補正は管の長さによらず一定とする。
\end{minipage}\hfill
\begin{minipage}[t]{0.45\linewidth}\vspace{0pt}\centering
\begin{tikzpicture}[line width=0.6pt,>={Stealth[length=1.8mm]}]
  \draw (0.6,0.45) -- (4.4,0.45);
  \draw (0.6,-0.45) -- (4.4,-0.45);
  \draw[very thick] (4.4,-0.6) -- (4.4,0.6);
  \draw[thick] (4.4,0) -- (4.9,0);
  \draw[<->] (0.6,-0.85) -- (4.4,-0.85) node[midway,below] {$L$};
  \draw[fill=gray!25] (-0.15,-0.2) rectangle (0.15,0.2);
  \draw[fill=gray!25] (0.15,0.2) -- (0.45,0.55) -- (0.45,-0.55) -- (0.15,-0.2) -- cycle;
  \node[above] at (0.0,0.25) {\small スピーカー};
  \node[above] at (2.4,0.55) {\small 管};
  \node[above] at (4.6,0.65) {\small 板};
\end{tikzpicture}\\[1pt]{\small 図3}
\end{minipage}
\begin{qlist}[resume*=sb]
\item 音の波長 $\lambda$ と、空気中の音速 $V$ を求めよ。\pts{8}
\item 開口端補正 $\Delta L$ を求めよ。\pts{6}
\item 3 回目に共鳴する管の長さ $L_3$ を求めよ。\pts{6}
\end{qlist}
\bigpart{Ⅲ}振動数 $720\unit{Hz}$ の音を出す音源が、静止した観測者に向かって一直線上を一定の速さ $20\unit{m/s}$ で近づき、観測者の前を通過したのち遠ざかった。音速は $340\unit{m/s}$ で、風はないものとする。
\begin{qlist}[resume*=sb]
\item 音源が近づく間に、観測者が聞く音の振動数 $f_1$ を求めよ。\pts{8}
\item 音源が遠ざかる間に、観測者が聞く音の振動数 $f_2$ を求めよ。\pts{8}
\item 音源を静止させ、観測者が速さ $20\unit{m/s}$ で音源に近づく場合に聞こえる振動数は、(8) の $f_1$ とわずかに異なる。その理由を、音の波長と、観測者から見た音の速さに着目して 3 行程度で述べよ。\pts{11}
\end{qlist}
\end{ascolorbox4A}

\clearpage
% ---------------------------- 大問3 ----------------------------
\begin{ascolorbox4A}[75点・30分]{大問3　電気}
図4の回路で、電池の起電力は $E=24\unit{V}$（内部抵抗は無視できる）、抵抗は $R_1=4.0\,\Omega$、$R_2=6.0\,\Omega$、$R_3=12\,\Omega$、コンデンサーの電気容量は $C=2.0\,\mu\mathrm{F}$ である。最初、コンデンサーに電荷は蓄えられておらず、スイッチ $\mathrm{S}_1$ は開き、$\mathrm{S}_2$ は閉じている。点 A、B は図の接続点である。
\begin{center}
\begin{circuitikz}[european,line width=0.6pt,scale=0.85,transform shape]
  \draw (0,0) to[battery1,l=$E$] (0,3.2);
  \draw (0,3.2) to[R,l=$R_1$] (3,3.2) -- (7.6,3.2);
  \draw (3.6,3.2) to[spst,l=$\mathrm{S}_1$] (3.6,1.9) to[C,l=$C$] (3.6,0);
  \draw (5.4,3.2) to[spst,l=$\mathrm{S}_2$] (5.4,1.9) to[R,l=$R_2$] (5.4,0);
  \draw (7.6,3.2) to[R,l=$R_3$] (7.6,0);
  \draw (0,0) -- (7.6,0);
  \fill (3.6,3.2) circle[radius=0.07] node[above left] {A};
  \fill (3.6,0) circle[radius=0.07] node[below left] {B};
  \fill (5.4,3.2) circle[radius=0.07];
  \fill (5.4,0) circle[radius=0.07];
\end{circuitikz}\\[1pt]{\small 図4}
\end{center}
\begin{qlist}[series=sc]
\item $\mathrm{S}_1$ を開いたままのとき、回路全体の合成抵抗と、電池を流れる電流 $I$ を求めよ。\pts{7}
\item このとき、$R_2$ と $R_3$ を流れる電流 $I_2$、$I_3$ をそれぞれ求めよ。\pts{7}
\item このとき、$R_1$ で消費される電力 $P_1$ と、$10\unit{s}$ 間に $R_1$ で発生するジュール熱 $H$ を求めよ。\pts{10}
\end{qlist}
次に、$\mathrm{S}_1$ を閉じて十分に時間がたった。
\begin{qlist}[resume*=sc]
\item コンデンサーの極板間の電圧 $V_C$、蓄えられた電気量 $Q$、静電エネルギー $U$ を求めよ。\pts{12}
\end{qlist}
その後、$\mathrm{S}_2$ を開いた。この時刻を $t=0$ とする。
\begin{qlist}[resume*=sc]
\item $t=0$ の直後に、$R_1$ を流れる電流 $I_1'$ と、コンデンサーに流れ込む電流 $I_C$ を求めよ。\pts{13}
\item 十分に時間がたったあとの、コンデンサーの電圧 $V_C'$ と蓄えられた電気量 $Q'$ を求めよ。また、$t=0$ から十分に時間がたつまでに、コンデンサーに新たに流れ込んだ電気量 $\Delta Q$ を求めよ。\pts{14}
\item $t\ge 0$ におけるコンデンサーの電圧 $V_C$ の時間変化の概形を、グラフに描け。$t=0$ のときの値と、十分に時間がたったあとの値を、縦軸に書き入れること。\pts{12}
\end{qlist}
\end{ascolorbox4A}

\clearpage
% ---------------------------- 大問4 ----------------------------
\begin{ascolorbox4A}[75点・30分]{大問4　力学（運動量・円運動・単振動）}
図5のように、水平でなめらかな床の上に、左端を壁に固定した軽いばね（ばね定数 $k$）がある。ばねの右端に質量 $m$ の小物体 A を押しつけ、ばねを自然の長さから $d$ だけ縮めて静かに放す。A はばねの自然の長さの位置 $\mathrm{O}$ でばねから離れ、床の上に静止している質量 $m$ の小物体 B に衝突して一体となった。一体となった物体は、床の点 P から、鉛直面内に置かれた半径 $R$ のなめらかな半円筒の内側面に沿って上り、最高点 Q から水平に飛び出して床に落下した。重力加速度の大きさを $g$ とし、床と半円筒の間にすきまはなく、物体の大きさは無視できるものとする。
\begin{center}
\begin{tikzpicture}[>={Stealth[length=2mm]},line width=0.6pt,x=1.0cm,y=1.0cm]
  \fill[pattern={Lines[angle=45,distance=2mm]}] (-0.3,-0.3) rectangle (9.6,0);
  \draw (-0.3,0) -- (9.6,0);
  \fill[pattern={Lines[angle=45,distance=2mm]}] (-0.3,0) rectangle (0,1.1);
  \draw (0,0) -- (0,1.1);
  \draw[decorate,decoration={coil,segment length=2.4mm,amplitude=1.6mm,aspect=0.5}] (0,0.38) -- (1.75,0.38);
  \draw[fill=gray!20] (1.75,0) rectangle (2.55,0.76);  \node at (2.15,0.38) {A};
  \draw[fill=gray!20] (4.6,0) rectangle (5.4,0.76);    \node at (5.0,0.38) {B};
  \node[above] at (2.15,0.76) {\small $m$}; \node[above] at (5.0,0.76) {\small $m$};
  \draw[densely dashed] (2.75,0) -- (2.75,1.1);  \node[above] at (2.75,1.1) {$\mathrm{O}$};
  \draw[<->] (1.75,1.0) -- (2.75,1.0) node[midway,above] {\small $d$};
  \draw[thick] (7.3,0) arc[start angle=-90,end angle=90,radius=1.5];
  \draw[densely dashed] (7.3,0) -- (7.3,3.0);
  \fill (7.3,0) circle[radius=0.06] node[below,fill=white,inner sep=1pt,yshift=-1.5mm] {P};
  \fill (7.3,3.0) circle[radius=0.06] node[above] {Q};
  \draw (7.3,1.5) -- ++(30:1.5) node[midway,above] {\small $R$};
  \fill (7.3,1.5) circle[radius=0.04];
  \draw[->,thick] (7.3,3.0) -- ++(-1.1,0) node[left] {$v_\mathrm{Q}$};
\end{tikzpicture}\\[1pt]{\small 図5}
\end{center}
\begin{qlist}[series=sd]
\item A がばねから離れる直前の速さ $v_0$ を、$k,m,d$ で表せ。\pts{5}
\item A を放してから、A がばねから離れるまでの時間 $t_1$ を、$k,m$ で表せ。\pts{7}
\item A と B が一体となった直後の速さ $v_1$ と、この衝突で失われた力学的エネルギー $\Delta E$ を、それぞれ $k,m,d$ で表せ。\pts{9}
\item 一体の物体が点 P を通過する瞬間に、半円筒から受ける垂直抗力の大きさ $N_\mathrm{P}$ を、$k,m,d,g,R$ で表せ。\pts{9}
\item 一体の物体が点 Q に達したときの速さ $v_\mathrm{Q}$ と、そのとき半円筒から受ける垂直抗力の大きさ $N_\mathrm{Q}$ を、$v_1,m,g,R$ で表せ。\pts{11}
\item 一体の物体が、半円筒の内側面から離れることなく点 Q に達するために、$d$ が満たすべき条件を、$k,m,g,R$ で表せ。\pts{11}
\item 点 Q から飛び出した物体が、床に落下するまでの時間と、落下点の点 P からの距離 $x$ を、$v_\mathrm{Q},g,R$ で表せ。\pts{8}
\item 落下点の点 P からの距離が $x=2\sqrt{3}\,R$ であったとする。このときの $d$ を、$k,m,g,R$ で表せ。\pts{15}
\end{qlist}
\end{ascolorbox4A}
\clearpage
% ==================================================================
%  問題編　セット2
% ==================================================================
\begin{simple}[120分・300点]{予想問題　セット2　注意事項}
\begin{itemize}[leftmargin=1.2\zw,itemsep=0pt,topsep=1pt,parsep=0pt]
\item 4 大問（各 75 点・目標 30 分）。大問は易しいものから難しいものへ並べてある。電卓・定規・コンパスは使えないものとし、図やグラフは徒手で描く。
\item 重力加速度の大きさは $g=9.8\unit{m/s^2}$ とする（文字 $g$ で答える設問は文字のままでよい）。
\item 数値で答える設問は、指定がなければ有効数字 2 桁とし、単位を必ず付ける（割り切れる場合は、そのままの値でよい）。$\sqrt{\phantom{x}}$、$\pi$、分数は計算せずに残してよい。
\item 空気の抵抗は無視する。答えだけでなく、考え方と途中の式も記すこと。
\end{itemize}
\end{simple}

% ---------------------------- 大問1 ----------------------------
\begin{ascolorbox4A}[75点・30分]{大問1　力学（斜面と円運動）}
\par\noindent
\begin{minipage}[t]{0.50\linewidth}\vspace{0pt}
\bpart{Ⅰ}図1のように、水平な床の上に、傾きの角 $\theta=30^\circ$ の固定された斜面がある。斜面の下端から、質量 $m$ の小物体を斜面に沿って上向きに初速 $v_0$ で打ち出した。物体と斜面の間の動摩擦係数は $\mu'=\dfrac{\sqrt{3}}{6}$、静止摩擦係数は $\mu=\dfrac{1}{2}$ である。重力加速度の大きさを $g$ とし、$\sin30^\circ=\dfrac{1}{2}$、$\cos30^\circ=\dfrac{\sqrt{3}}{2}$ を用いてよい。
\end{minipage}\hfill
\begin{minipage}[t]{0.47\linewidth}\vspace{0pt}\centering
\begin{tikzpicture}[>={Stealth[length=2mm]},line width=0.6pt,x=0.8cm,y=0.8cm]
  \def\ang{30}
  \fill[pattern={Lines[angle=45,distance=2mm]}] (0,0) -- (6,0) -- (6,-0.25) -- (0,-0.25) -- cycle;
  \draw (0,0) -- (6,0) -- (6,{6*tan(\ang)}) -- cycle;
  \draw (1.3,0) arc[start angle=0,end angle=\ang,radius=1.3];
  \node at (1.75,0.3) {\small $\theta$};
  \begin{scope}[rotate=\ang]
    \draw[fill=gray!20] (0.5,0) rectangle (1.4,0.7);
    \node at (0.95,0.35) {$m$};
    \draw[->,very thick] (1.55,0.35) -- (2.75,0.35) node[above] {$v_0$};
  \end{scope}
\end{tikzpicture}\\[1pt]{\small 図1}
\end{minipage}
\begin{qlist}[series=ta]
\item 物体が斜面を上向きに滑っているとき、物体にはたらく力をすべて、図1の物体に矢印で描き入れ、それぞれの名称を書け。\pts{6}
\item 上向きに滑っている間の加速度の大きさ $a_1$ を、$g$ で表せ。\pts{7}
\item 物体が最高点に達するまでに斜面に沿って進む距離 $s$ を、$v_0,g$ で表せ。また、$v_0=7.0\unit{m/s}$ のときの $s$ を、有効数字 2 桁で求めよ。\pts{6}
\item 最高点に達した物体は、そのまま静止し続けるか、それとも斜面を滑り降りるか。理由とともに答えよ。\pts{7}
\item 物体が滑り降りる場合、斜面の下端にもどってきたときの速さ $v$ を、$v_0$ で表せ。\pts{9}
\end{qlist}
\par\noindent
\begin{minipage}[t]{0.55\linewidth}\vspace{0pt}
\bpart{Ⅱ}質量 $m$ の小球を長さ $L$ の軽い糸の一端につけ、他端を天井の点 $\mathrm{O}$ に固定する。図2のように、小球が水平面内で等速円運動をしているとき、糸は鉛直線と角 $\theta$ をなしていた。円運動の半径を $r$、重力加速度の大きさを $g$ とする。
\end{minipage}\hfill
\begin{minipage}[t]{0.42\linewidth}\vspace{0pt}\centering
\begin{tikzpicture}[>={Stealth[length=2mm]},line width=0.6pt,x=0.85cm,y=0.85cm]
  \fill[pattern={Lines[angle=45,distance=2mm]}] (-2.0,0) rectangle (2.0,0.25);
  \draw (-2.0,0) -- (2.0,0);
  \def\h{2.0}\def\r{1.5}
  \coordinate (O) at (0,0);
  \coordinate (C) at (0,-\h);
  \coordinate (B) at (\r,-\h);
  \draw[densely dashed] (O) -- (C);
  \draw[densely dashed] (C) ellipse[x radius=\r,y radius=0.4];
  \draw (O) -- (B) node[midway,above right] {$L$};
  \draw[<->] (C) -- (B) node[midway,above] {\small $r$};
  \draw[fill=gray!35] (B) circle[radius=0.17];
  \node[right] at ($(B)+(0.2,0)$) {$m$};
  \draw (0,-0.85) arc[start angle=-90,end angle={-90+atan(\r/\h)},radius=0.85];
  \node at (0.28,-1.2) {\small $\theta$};
  \node[above] at (0,0.25) {$\mathrm{O}$};
\end{tikzpicture}\\[1pt]{\small 図2}
\end{minipage}
\begin{qlist}[resume*=ta]
\item 糸の張力の大きさ $S$ と、小球の速さ $v$ を、$m,g,L,\theta$ で表せ。\pts{11}
\item 円運動の周期 $T$ を、$L,g,\theta$ で表せ。\pts{10}
\item $\theta=60^\circ$、$L=4.9\unit{m}$、$m=2.0\unit{kg}$ のとき、周期 $T$ と張力の大きさ $S$ を、有効数字 2 桁で求めよ。\pts{9}
\item $L$ を一定にしたまま小球の回転を速くしていくと、$\theta$ は大きくなる。このとき周期 $T$ はどうなるか。理由とともに述べよ。\pts{10}
\end{qlist}
\end{ascolorbox4A}

\clearpage
% ---------------------------- 大問2 ----------------------------
\begin{ascolorbox4A}[75点・30分]{大問2　電磁気（磁場中の荷電粒子・電磁誘導）}
\par\noindent
\begin{minipage}[t]{0.52\linewidth}\vspace{0pt}
\bpart{Ⅰ}図3のように、質量 $m$、電気量 $q\,(>0)$ の粒子を、電圧 $V$ で静止の状態から加速し、点 $\mathrm{O}$ から一様な磁場の領域へ、境界に垂直に入射させる。磁束密度の大きさは $B$ で、向きは紙面に垂直で手前向き（図の $\odot$）である。重力と空気の抵抗は無視する。
\end{minipage}\hfill
\begin{minipage}[t]{0.45\linewidth}\vspace{0pt}\centering
\begin{tikzpicture}[>={Stealth[length=2mm]},line width=0.6pt,x=0.78cm,y=0.78cm]
  \fill[gray!8] (0,-3.7) rectangle (3.4,0.8);
  \draw (0,-3.7) -- (0,0.8);
  \foreach \x in {0.7,1.5,2.3,3.0} \foreach \y in {0.4,-0.4,-1.2,-2.0,-2.8,-3.4}
    {\draw (\x,\y) circle[radius=0.07]; \fill (\x,\y) circle[radius=0.018];}
  \draw[very thick] (-2.7,-0.5) -- (-2.7,0.5);
  \draw[very thick] (-1.8,-0.5) -- (-1.8,-0.12) (-1.8,0.12) -- (-1.8,0.5);
  \node[above] at (-2.7,0.5) {\small $+$}; \node[above] at (-1.8,0.5) {\small $-$};
  \draw (-2.7,-0.5) -- (-2.7,-1.4) -- (-2.35,-1.4);
  \draw (-2.35,-1.65) -- (-2.35,-1.15);
  \draw (-2.15,-1.5) -- (-2.15,-1.3);
  \draw (-2.15,-1.4) -- (-1.8,-1.4) -- (-1.8,-0.5);
  \node[below] at (-2.25,-1.65) {$V$};
  \fill[gray!60] (-2.55,0) circle[radius=0.1];
  \node[below] at (-2.25,-0.1) {\small $q,m$};
  \draw[->,thick] (-2.55,0) -- (-0.4,0);
  \draw[->,thick] (0,0) arc[start angle=90,end angle=-90,radius=1.5];
  \draw[<->] (0,-1.5) -- (1.5,-1.5);
  \node[above] at (0.8,-1.5) {\small $r$};
  \fill (0,0) circle[radius=0.06] node[above left] {O};
  \fill (0,-3.0) circle[radius=0.06] node[below left] {P};
  \node at (2.45,0.78) {\small $B$};
\end{tikzpicture}\\[1pt]{\small 図3}
\end{minipage}
\begin{qlist}[series=tb]
\item 点 $\mathrm{O}$ に入射するときの粒子の速さ $v$ を、$q,m,V$ で表せ。\pts{6}
\item 磁場中の円運動の半径 $r$ と周期 $T$ を求めよ。ただし $r$ は $q,m,V,B$ で、$T$ は $q,m,B$ で表せ。\pts{10}
\item 粒子は境界上の点 $\mathrm{P}$ から磁場の外へ出た。$\mathrm{OP}$ の長さを $r$ で表せ。\pts{4}
\item 質量 $2m$、電気量 $q$ の粒子を、同じ電圧 $V$ で加速して同じように入射させた。このときの円運動の半径は、(2) の $r$ の何倍か。\pts{6}
\item 磁場の中で、粒子の速さが変わらないのはなぜか。1 文で述べよ。\pts{6}
\end{qlist}
\par\noindent
\begin{minipage}[t]{0.50\linewidth}\vspace{0pt}
\bpart{Ⅱ}図4は、水平面内に置かれた間隔 $\ell=0.40\unit{m}$ の平行な導体レールを真上から見た図である。レールの左端は抵抗 $R=2.0\,\Omega$ でつながれ、レール上に導体棒 PQ がレールと直角に置かれている。鉛直上向き（紙面に垂直で手前向き）に、磁束密度 $B=0.50\unit{T}$ の一様な磁場がある。棒 PQ を右向きに一定の速さ $v=5.0\unit{m/s}$ で動かす。レールと棒の電気抵抗、および摩擦は無視できる。
\end{minipage}\hfill
\begin{minipage}[t]{0.47\linewidth}\vspace{0pt}\centering
\begin{tikzpicture}[>={Stealth[length=2mm]},line width=0.6pt,x=0.8cm,y=0.8cm]
  \draw (0,0) -- (5.8,0);
  \draw (0,2) -- (5.8,2);
  \draw (0,0) -- (0,0.7) (0,1.3) -- (0,2);
  \draw[fill=white] (-0.18,0.7) rectangle (0.18,1.3);
  \node[left] at (-0.25,1) {$R$};
  \foreach \x in {1.0,1.9,2.8,3.7,4.6,5.4} \foreach \y in {0.5,1.0,1.5}
    {\draw (\x,\y) circle[radius=0.07]; \fill (\x,\y) circle[radius=0.018];}
  \draw[very thick] (3.2,-0.3) -- (3.2,2.3);
  \node[above] at (3.2,2.3) {P}; \node[below] at (3.2,-0.3) {Q};
  \draw[->,very thick] (3.35,1.0) -- (4.3,1.0) node[above] {$v$};
  \draw[<->] (5.2,0) -- (5.2,2) node[midway,right] {$\ell$};
  \node at (0.55,1.78) {\small $B$};
\end{tikzpicture}\\[1pt]{\small 図4}
\end{minipage}
\begin{qlist}[resume*=tb]
\item 棒 PQ に生じる誘導起電力の大きさ $V_0$ を求めよ。\pts{6}
\item 抵抗 $R$ を流れる電流の大きさ $I$ を求めよ。また、棒の中を流れる電流の向き（P→Q か Q→P か）と、P、Q のうち電位が高い方を答えよ。\pts{10}
\item 棒を一定の速さで動かし続けるために必要な外力の大きさ $F$ と、その外力の仕事率 $P$ を求めよ。\pts{10}
\item 抵抗で単位時間あたりに発生するジュール熱 $P_R$ を求め、(8) の $P$ との関係を、エネルギーの観点から述べよ。\pts{7}
\item 棒を押す外力を急に取り去った。その後、棒の速さはどのように変化するか。理由とともに述べよ。\pts{10}
\end{qlist}
\end{ascolorbox4A}

\clearpage
% ---------------------------- 大問3 ----------------------------
\begin{ascolorbox4A}[75点・30分]{大問3　波（光の干渉）}
\bpart{Ⅰ}図5のように、波長 $\lambda=5.0\times10^{-7}\unit{m}$ の単色光を、間隔 $d=0.20\unit{mm}$ の 2 つのスリット $\mathrm{S}_1$、$\mathrm{S}_2$ に垂直に当て、スリットから $L=1.0\unit{m}$ はなれたスクリーンに干渉縞をつくる。スクリーン上で、2 つのスリットから等距離にある点を $\mathrm{O}$ とし、$\mathrm{O}$ から $\mathrm{S}_1$ の側へ測った距離を $y$ とする。$d$ と $|y|$ は $L$ に比べて十分に小さい。装置は空気中にあり、空気の屈折率は $1$ とする。
\begin{center}
\begin{tikzpicture}[line width=0.6pt,>={Stealth[length=1.8mm]},x=1cm,y=1.05cm]
  \coordinate (S1) at (0,0.5); \coordinate (S2) at (0,-0.5);
  \coordinate (P) at (8,1.5);  \coordinate (O) at (8,0);
  \draw[very thick] (0,-1.6) -- (0,-0.58) (0,-0.42) -- (0,0.42) (0,0.58) -- (0,1.6);
  \draw[very thick] (8,-1.6) -- (8,2.3);
  \draw[densely dashed] (0,0) -- (O);
  \draw (S1) -- (P); \draw (S2) -- (P);
  \pgfmathsetmacro{\rone}{veclen(8,1.0)}
  \pgfmathsetmacro{\rtwo}{veclen(8,2.0)}
  \pgfmathsetmacro{\ratio}{\rone/\rtwo}
  \coordinate (H) at ($(P)!\ratio!(S2)$);
  \draw[line width=1.6pt] (H) -- (S2);
  \node[below right,yshift=-1mm] at ($(S2)!0.5!(H)$) {\small 経路差};
  \draw[<->] (-1.0,-0.5) -- (-1.0,0.5) node[midway,left] {$d$};
  \draw[<->] (0,-1.95) -- (8,-1.95) node[midway,below] {$L$};
  \draw[<->] (8.45,0) -- (8.45,1.5) node[midway,right] {$y$};
  \fill (S1) circle[radius=0.04] node[left,xshift=-0.05cm] {$\mathrm{S}_1$};
  \fill (S2) circle[radius=0.04] node[left,xshift=-0.05cm] {$\mathrm{S}_2$};
  \fill (P) circle[radius=0.05] node[left] {$\mathrm{P}$};
  \fill (O) circle[radius=0.05] node[left] {$\mathrm{O}$};
\end{tikzpicture}\\[1pt]{\small 図5}
\end{center}
\begin{qlist}[series=tc]
\item 屈折率 $n$ の媒質中を進む光の波長・速さ・振動数は、それぞれ真空中の何倍か。\pts{6}
\item スクリーン上の点 $\mathrm{P}$（座標 $y$）に届く 2 つの光の経路差 $|\mathrm{S}_2\mathrm{P}-\mathrm{S}_1\mathrm{P}|$ を、$d,y,L$ で表せ。\pts{8}
\item 明線ができる位置 $y_m$（$m=0,\pm1,\pm2,\dots$）と、隣り合う明線の間隔 $\Delta y$ を求めよ。$\Delta y$ は数値でも答えよ。\pts{10}
\end{qlist}
\par\noindent
\begin{minipage}[t]{0.56\linewidth}\vspace{0pt}
次に、図6のように、$\mathrm{S}_1$ の直後に、屈折率 $n=1.5$、厚さ $t=5.0\times10^{-6}\unit{m}$ の薄いガラス板を、光の進行方向に垂直に置いた。光がこのガラス板を通過すると、光路長（屈折率と進んだ距離の積）は、ガラス板がない場合に比べて $(n-1)t$ だけ長くなる。
\end{minipage}\hfill
\begin{minipage}[t]{0.40\linewidth}\vspace{0pt}\centering
\begin{tikzpicture}[line width=0.6pt,>={Stealth[length=1.8mm]},x=1.1cm,y=1.1cm]
  \draw[very thick] (0,-1.5) -- (0,-0.95) (0,-0.65) -- (0,0.65) (0,0.95) -- (0,1.5);
  \draw[fill=gray!25] (0.2,0.5) rectangle (0.75,1.1);
  \node at (0.47,1.45) {\small ガラス板};
  \draw[<->] (0.2,0.3) -- (0.75,0.3);
  \node[below] at (0.47,0.3) {\small $t$};
  \draw[->] (-1.1,0.8) -- (-0.1,0.8); \draw[->] (-1.1,-0.8) -- (-0.1,-0.8);
  \draw[->] (0.75,0.8) -- (1.8,0.8);  \draw[->] (0.1,-0.8) -- (1.8,-0.8);
  \node[left] at (-1.1,0.8) {$\mathrm{S}_1$}; \node[left] at (-1.1,-0.8) {$\mathrm{S}_2$};
  \node at (0.47,0.8) {\small $n$};
\end{tikzpicture}\\[1pt]{\small 図6}
\end{minipage}
\begin{qlist}[resume*=tc]
\item このとき、点 $\mathrm{P}$ が明線となる条件を、$d,y,L,t,n,\lambda,m$ を用いて表せ（$m$ は整数）。\pts{8}
\item 中央の明線（$m=0$ の明線）の位置 $y_0$ を求めよ。また、その明線は、ガラス板を置く前の明線の間隔 $\Delta y$ の何倍分、$\mathrm{S}_1$ 側と $\mathrm{S}_2$ 側のどちらへずれたか。\pts{10}
\end{qlist}
\bpart{Ⅱ}1 mm あたり 500 本の溝がある回折格子に、同じ単色光（波長 $5.0\times10^{-7}\unit{m}$）を格子面に垂直に当て、格子の後方に置いたスクリーンに明線をつくる。
\begin{qlist}[resume*=tc]
\item 格子定数 $d'$ を求めよ。また、2 次の明線の回折角 $\theta_2$ を求めよ。\pts{10}
\item スクリーン上に観測される明線は、中央の明線をふくめて全部で何本か。理由とともに答えよ。\pts{10}
\item 単色光のかわりに白色光を当てると、1 次の明線は、中央に近い側が紫色、遠い側が赤色になる。その理由を述べよ。\pts{13}
\end{qlist}
\end{ascolorbox4A}

\clearpage
% ---------------------------- 大問4 ----------------------------
\begin{ascolorbox4A}[75点・30分]{大問4　熱と気体（ピストン付きシリンダー）}
\par\noindent
\begin{minipage}[t]{0.58\linewidth}\vspace{0pt}
鉛直に立てた断熱材のシリンダーに、質量 $M$ のなめらかに動くピストンで、物質量 $n$ の単原子分子理想気体が閉じ込められている。ピストンの断面積は $S$、外気の圧力は $p_0$、気体定数は $R$、重力加速度の大きさは $g$ である。シリンダー内にはヒーターがあり、気体を加熱できる（ヒーターの体積は無視する）。また、外部の冷却装置で気体を冷やすこともできる。図7の状態 A では、気体の温度は $T_0$、ピストンはシリンダーの底から高さ $h_0$ の位置で静止している。単原子分子理想気体の内部エネルギーは $U=\dfrac{3}{2}nRT$ である。
\end{minipage}\hfill
\begin{minipage}[t]{0.39\linewidth}\vspace{0pt}\centering
\begin{tikzpicture}[>={Stealth[length=2mm]},line width=0.6pt,x=0.9cm,y=0.9cm]
  \fill[blue!7] (0.03,0.03) rectangle (2.37,1.6);
  \fill[pattern={Lines[angle=45,distance=2mm]}] (-0.22,-0.22) rectangle (2.62,0);
  \fill[pattern={Lines[angle=45,distance=2mm]}] (-0.22,0) rectangle (0,3.3);
  \fill[pattern={Lines[angle=45,distance=2mm]}] (2.4,0) rectangle (2.62,3.3);
  \draw[very thick] (0,3.3) -- (0,0) -- (2.4,0) -- (2.4,3.3);
  \draw[fill=gray!35] (0.03,1.6) rectangle (2.37,1.9);
  \node at (1.2,1.75) {\small $M$};
  \node at (1.2,1.0) {\small 気体};
  \draw[decorate,decoration={zigzag,segment length=2mm,amplitude=0.7mm}] (0.4,0.3) -- (2.0,0.3);
  \node at (1.2,0.55) {\tiny ヒーター};
  \foreach \x in {0.5,1.2,1.9} \draw[->] (\x,2.9) -- (\x,2.0);
  \node[above] at (1.2,2.9) {$p_0$};
  \draw[<->] (2.85,0) -- (2.85,1.6) node[midway,right] {$h_0$};
  \draw[densely dashed] (2.4,1.6) -- (2.95,1.6);
\end{tikzpicture}\\[1pt]{\small 図7}
\end{minipage}
\begin{qlist}[series=td]
\item 状態 A の気体の圧力 $p_1$ を、$p_0,M,g,S$ で表せ。\pts{5}
\item $h_0$ を、$n,R,T_0,S,p_0,M,g$ で表せ。\pts{8}
\end{qlist}
ヒーターで気体をゆっくり加熱したところ、ピストンはゆっくり上昇し、高さが $2h_0$ になった（状態 B）。
\begin{qlist}[resume*=td]
\item 状態 B の気体の温度 $T_\mathrm{B}$ を、$T_0$ で表せ。\pts{6}
\item A→B の間に気体が吸収した熱量 $Q_\mathrm{AB}$ を、$n,R,T_0$ で表せ。\pts{10}
\item A→B の間に気体がした仕事 $W$ のうち、(i) ピストンの重力による位置エネルギーの増加分と、(ii) 大気に対してした仕事を、それぞれ $M,g,p_0,S,h_0$ で表せ。\pts{8}
\item A→B で気体の温度は $2$ 倍になったが、圧力は変化しなかった。この理由を、気体分子の速さ、単位体積あたりの分子数、壁に衝突する分子のようすに着目して、3 行程度で述べよ。\pts{12}
\end{qlist}
次に、ピストンをストッパーで固定し、冷却装置で気体を冷やして、温度を $T_0$ にもどした（状態 C）。
\begin{qlist}[resume*=td]
\item 状態 C の圧力 $p_\mathrm{C}$ を $p_1$ で表せ。また、B→C で気体が放出した熱量 $Q_\mathrm{BC}$ を、$n,R,T_0$ で表せ。\pts{10}
\item A→B→C の変化を、縦軸を $p$、横軸を $V$ とするグラフに描け。状態 A, B, C の座標（$p_1$ や $V_0=Sh_0$ を用いてよい）を書き入れること。\pts{8}
\item A→B→C の全体で気体が吸収した正味の熱量 $Q_\mathrm{net}$ を、$n,R,T_0$ で表せ。また、$Q_\mathrm{net}$ が気体のした正味の仕事 $W_\mathrm{net}$ に等しくなる理由を述べよ。\pts{8}
\end{qlist}
\end{ascolorbox4A}
\clearpage
% ==================================================================
%  問題編　セット3（追加：力学と電気を中心にした、やや難しめのセット）
% ==================================================================
\begin{simple}[120分・300点]{予想問題　セット3（追加）　注意事項}
\begin{itemize}[leftmargin=1.2\zw,itemsep=0pt,topsep=1pt,parsep=0pt]
\item セット1・2 より力学と電気の比重を高め、難度もやや上げた追加セット（標準〜やや応用）。4 大問（各 75 点・目標 30 分）で、後ろの大問ほど難しい。
\item 重力加速度の大きさは $g=9.8\unit{m/s^2}$ とする（文字 $g$ で答える設問は文字のままでよい）。
\item 数値で答える設問は、指定がなければ有効数字 2 桁とし、単位を必ず付ける（割り切れる場合は、そのままの値でよい）。$\sqrt{\phantom{x}}$、$\pi$、分数は計算せずに残してよい。
\item 空気の抵抗は無視する。答えだけでなく、考え方と途中の式も記すこと。
\end{itemize}
\end{simple}

% ---------------------------- 大問1 ----------------------------
\begin{ascolorbox4A}[75点・30分]{大問1　電気（コンデンサーと電荷保存）}
図1の回路で、電池の起電力は $E=12\unit{V}$、コンデンサーの電気容量は $C_1=3.0\,\mu\mathrm{F}$、$C_2=6.0\,\mu\mathrm{F}$ である。はじめ、スイッチ $\mathrm{S}_1$、$\mathrm{S}_2$ は開いており、2 つのコンデンサーに電荷は蓄えられていない。電池と導線の内部抵抗以外の抵抗は考えず、スイッチを操作したあとは、十分に時間がたってから次の操作に移るものとする。
\begin{center}
\begin{circuitikz}[european,line width=0.6pt,scale=0.85,transform shape]
  \draw (0,0) to[battery1,l=$E$] (0,3);
  \draw (0,3) to[spst,l=$\mathrm{S}_1$] (2.4,3) -- (6.0,3);
  \draw (3.0,3) to[C,l=$C_1$] (3.0,0);
  \draw (6.0,3) to[spst,l=$\mathrm{S}_2$] (6.0,1.9) to[C,l=$C_2$] (6.0,0);
  \draw (0,0) -- (6.0,0);
  \fill (3.0,3) circle[radius=0.07];
  \fill (3.0,0) circle[radius=0.07];
\end{circuitikz}\\[1pt]{\small 図1}
\end{center}
\begin{qlist}[series=ua]
\item $\mathrm{S}_1$ だけを閉じた。$C_1$ に蓄えられる電気量 $Q_1$ と、静電エネルギー $U_1$ を求めよ。\pts{8}
\item 次に、$\mathrm{S}_1$ を開き、その後 $\mathrm{S}_2$ を閉じた。2 つのコンデンサーの電圧 $V'$ と、$C_1$、$C_2$ に蓄えられる電気量 $Q_1'$、$Q_2'$ を求めよ。\pts{12}
\item (1) の状態から (2) の状態になる間に失われた静電エネルギー $\Delta U$ を求めよ。また、失われたエネルギーは何に変わったか。\pts{8}
\item 続いて、$\mathrm{S}_2$ を閉じたまま $\mathrm{S}_1$ を閉じた。この間に電池の中を移動した電気量 $\Delta Q$ と、電池が電荷にした仕事 $W$ を求めよ。\pts{12}
\item (4) の間に、コンデンサーに蓄えられた静電エネルギーの増加量 $\Delta U'$ と、回路で発生したジュール熱 $H$ を求めよ。\pts{10}
\item 次に、$\mathrm{S}_1$ を開き、$\mathrm{S}_2$ は閉じたまま、$C_2$ の極板間をすきまなく、比誘電率 $\varepsilon_\mathrm{r}=2.0$ の誘電体で満たした。2 つのコンデンサーの電圧 $V''$ を求めよ。\pts{12}
\item (6) の操作の前後で、静電エネルギーの合計はどのように変化したか。変化量を求め、その理由を、誘電体にはたらく静電気力の仕事に着目して述べよ。\pts{13}
\end{qlist}
\end{ascolorbox4A}

\clearpage
% ---------------------------- 大問2 ----------------------------
\begin{ascolorbox4A}[75点・30分]{大問2　力学（板の上の物体と運動量保存）}
図2のように、なめらかな水平面上に質量 $2m$ の板 B を置き、板の左端に質量 $m$ の小物体 A をのせる。右向きを正とする。A に初速 $v_0$ を与えると、A は板の上を右向きに滑り出した。A と B の間の動摩擦係数を $\mu'$、重力加速度の大きさを $g$ とし、A の大きさは無視できる。板の右には固定された壁があり、板は壁に弾性衝突する（衝突の時間は極めて短く、衝突の間に A にはたらく摩擦力の力積は無視できる）。壁は十分遠く、板が壁に達するまでに、A と B の速さは等しくなっているものとする。
\begin{center}
\begin{tikzpicture}[>={Stealth[length=2mm]},line width=0.6pt,x=1.0cm,y=1.0cm]
  \fill[pattern={Lines[angle=45,distance=2mm]}] (-0.3,-0.3) rectangle (9.6,0);
  \draw (-0.3,0) -- (9.6,0);
  \fill[pattern={Lines[angle=45,distance=2mm]}] (9.0,0) rectangle (9.3,1.9);
  \draw (9.0,0) -- (9.0,1.9);
  \draw[fill=gray!12] (0.4,0) rectangle (6.4,0.45);
  \node at (4.2,0.23) {\small $\mathrm{B}$（$2m$）};
  \draw[fill=gray!35] (0.5,0.45) rectangle (1.3,1.15);
  \node at (0.9,0.8) {A};
  \node[above] at (0.9,1.15) {\small $m$};
  \draw[->,very thick] (1.45,0.8) -- (2.75,0.8) node[above] {$v_0$};
  \node[above] at (4.9,0.47) {\small 動摩擦係数 $\mu'$};
  \draw[<->] (0.4,-0.6) -- (6.4,-0.6) node[midway,below] {\small 板の長さ $\ell_0$};
\end{tikzpicture}\\[1pt]{\small 図2}
\end{center}
\begin{qlist}[series=ub]
\item A と B が滑っている間の、A の加速度の大きさ $a$ と、B の加速度の大きさ $b$ を、$\mu',g$ で表せ。\pts{6}
\item A と B の速さが等しくなるまでの時間 $t_1$ と、そのときの共通の速さ $v_1$ を求めよ。\pts{10}
\item (2) までの間に、A が B の上を滑った距離 $\ell_1$ を、$v_0,\mu',g$ で表せ。また、$v_0=7.0\unit{m/s}$、$\mu'=0.50$ のときの $\ell_1$ を求めよ。\pts{12}
\item 時刻 $t=0$ から $t=2t_1$ までの、A の速度と B の速度の時間変化を、同じ座標軸にグラフで描け（板は $t=2t_1$ までは壁に達しないとする）。$t_1$、$v_0$、$v_1$ を軸に書き入れること。\pts{10}
\item 板が壁に弾性衝突した直後の、A の速度と B の速度を求めよ。その後、A と B の速度が等しくなったときの共通の速度 $V$ を、向きも含めて求めよ。\pts{15}
\item 壁に衝突してから A と B の速度が等しくなるまでに、A が B の上をさらに滑る距離 $\ell_2$ を求めよ。\pts{12}
\item A が板の右端から落ちないための、板の長さ $\ell_0$ の条件を求めよ。\pts{10}
\end{qlist}
\end{ascolorbox4A}

\clearpage
% ---------------------------- 大問3 ----------------------------
\begin{ascolorbox4A}[75点・30分]{大問3　電磁気（速度選択器と質量分析）}
図3は、速度選択器と質量分析器を組み合わせた装置である。質量 $m$、電気量 $q\,(>0)$ の粒子が、スリット $\mathrm{S}_1$ から右向きに、平行な 2 枚の極板の間へ入る。極板の間隔は $d$、電圧は $V$ で、上の極板が高電位である。この領域には、紙面に垂直な向きに大きさ $B_1$ の一様な磁場もあり、特定の速さ $v$ の粒子だけが直進して、スリット $\mathrm{S}_2$（点 $\mathrm{O}$）を通過する。点 $\mathrm{O}$ の右側には、紙面に垂直で手前向きの、大きさ $B_2$ の一様な磁場があり、粒子は半円を描いて、境界上の点 $\mathrm{P}$ に達する。重力は無視する。
\begin{center}
\begin{tikzpicture}[>={Stealth[length=2mm]},line width=0.6pt,x=0.82cm,y=0.82cm]
  \fill[gray!8] (0,-3.9) rectangle (3.4,1.4);
  \foreach \x in {0.7,1.5,2.3,3.0} \foreach \y in {1.0,0.2,-0.6,-1.4,-2.2,-3.0}
    {\draw (\x,\y) circle[radius=0.07]; \fill (\x,\y) circle[radius=0.018];}
  \draw[very thick] (0,1.4) -- (0,0.12) (0,-0.12) -- (0,-3.9);
  \draw[very thick] (-4.4,1.4) -- (-4.4,0.12) (-4.4,-0.12) -- (-4.4,-1.4);
  \draw[very thick] (-4.0,0.9) -- (-0.5,0.9);
  \draw[very thick] (-4.0,-0.9) -- (-0.5,-0.9);
  \node[above] at (-2.2,0.9) {$+$}; \node[below] at (-2.2,-0.9) {$-$};
  \foreach \x in {-3.4,-2.4,-1.4} \draw[->] (\x,0.8) -- (\x,0.35);
  \node[right] at (-1.15,0.58) {\small $E$};
  \foreach \x in {-3.6,-2.8,-2.0,-1.2}
    {\draw (\x,-0.5) circle[radius=0.1]; \draw (\x-0.07,-0.57) -- (\x+0.07,-0.43); \draw (\x-0.07,-0.43) -- (\x+0.07,-0.57);}
  \node[right] at (-0.95,-0.5) {\small $B_1$};
  \draw[<->] (-4.15,-0.9) -- (-4.15,0.9) node[midway,left] {$d$};
  \draw[->,thick] (-4.4,0) -- (-0.4,0);
  \draw[->,thick] (0,0) arc[start angle=90,end angle=-90,radius=1.5];
  \draw[<->] (0,-1.5) -- (1.5,-1.5); \node[above] at (0.75,-1.5) {\small $r$};
  \fill (0,0) circle[radius=0.06]; \node[above left] at (0,0.05) {O};
  \fill (0,-3.0) circle[radius=0.06]; \node[below left] at (0,-3.0) {P};
  \node[left] at (-4.5,0) {$\mathrm{S}_1$}; \node at (2.7,1.15) {\small $B_2$};
\end{tikzpicture}\\[1pt]{\small 図3}
\end{center}
\begin{qlist}[series=uc]
\item 速度選択器の中で粒子が直進するとき、磁場 $B_1$ の向きは、紙面に垂直で手前向きか奥向きか。また、直進する粒子の速さ $v$ を、$V,d,B_1$ で表せ。\pts{8}
\item $V=40\unit{V}$、$d=2.0\times10^{-2}\unit{m}$、$B_1=0.20\unit{T}$ のとき、$v$ を求めよ。\pts{6}
\item 点 $\mathrm{O}$ の右側での円運動の半径 $r$ を、$q,m,B_2,v$ で表せ。また、$\mathrm{OP}$ の長さから粒子の質量 $m$ を求める式を、$q,B_1,B_2,V,d,\mathrm{OP}$ で表せ。\pts{12}
\item $q=1.6\times10^{-19}\unit{C}$、$B_2=0.10\unit{T}$ とする。質量 $m=1.6\times10^{-26}\unit{kg}$ の粒子が描く円の半径 $r$ を求めよ。\pts{6}
\item 同じ電気量で質量 $1.8\times10^{-26}\unit{kg}$ の粒子（同位体）も、速度選択器を通過して入射した。この粒子が達する点を $\mathrm{P}'$ とする。$\mathrm{PP}'$ の長さを求めよ。\pts{12}
\item 速さが $v$ より大きい粒子は、スリット $\mathrm{S}_2$ を通過できない。その理由を述べよ。\pts{10}
\item 質量 $m=1.6\times10^{-26}\unit{kg}$ の粒子が、点 $\mathrm{O}$ から点 $\mathrm{P}$ に達するまでに要する時間 $t$ を求めよ。\pts{9}
\item この粒子は、スリット $\mathrm{S}_1$ に入る前に、静止の状態から電位差 $V_0$ で加速された。$V_0$ を求めよ。\pts{12}
\end{qlist}
\end{ascolorbox4A}

\clearpage
% ---------------------------- 大問4 ----------------------------
\begin{ascolorbox4A}[75点・30分]{大問4　力学（鉛直ばねと重ねた物体の単振動）}
\par\noindent
\begin{minipage}[t]{0.62\linewidth}\vspace{0pt}
図4のように、鉛直に立てた、ばね定数 $k=98\unit{N/m}$ の軽いばねの上端に、質量 $M=0.50\unit{kg}$ の薄い板を取り付け、板の上に質量 $m=1.5\unit{kg}$ の小物体 B をのせる。板と B は鉛直方向にだけ運動する。板と B がつり合って静止している位置を原点 $\mathrm{O}$ とし、鉛直上向きに $x$ 軸をとる。重力加速度の大きさは $g=9.8\unit{m/s^2}$ とし、空気の抵抗は無視できる。
\end{minipage}\hfill
\begin{minipage}[t]{0.33\linewidth}\vspace{0pt}\centering
\begin{tikzpicture}[>={Stealth[length=2mm]},line width=0.6pt,x=0.9cm,y=0.9cm]
  \fill[pattern={Lines[angle=45,distance=2mm]}] (-1.0,-0.25) rectangle (1.0,0);
  \draw (-1.0,0) -- (1.0,0);
  \draw[decorate,decoration={coil,segment length=2.2mm,amplitude=2.0mm,aspect=0.5}] (0,0) -- (0,2.3);
  \node[left] at (-0.35,1.15) {$k$};
  \draw[fill=gray!35] (-0.85,2.3) rectangle (0.85,2.45);
  \draw[fill=gray!12] (-0.45,2.45) rectangle (0.45,3.15);
  \node at (0,2.8) {B};
  \node[left] at (-0.9,2.37) {\small 板 $M$};
  \node[right] at (0.5,2.85) {\small $m$};
  \draw[densely dashed] (0.9,2.45) -- (1.6,2.45);
  \node[right] at (1.6,2.45) {\small $\mathrm{O}$};
  \draw[->,thick] (1.25,2.45) -- (1.25,3.5) node[above] {$x$};
\end{tikzpicture}\\[1pt]{\small 図4}
\end{minipage}
\begin{qlist}[series=ud]
\item つり合いの状態での、ばねの縮み $x_0$ を求めよ。\pts{8}
\item 板と B を、つり合いの位置から $d=0.25\unit{m}$ だけ押し下げて静かに放した。板と B が離れずに運動している間、その運動が単振動になる理由を述べ、周期 $T$ を求めよ。\pts{12}
\item 放した直後に、B が板から受ける垂直抗力の大きさ $N_0$ を求めよ。\pts{10}
\item B が板から離れるのは、ばねの長さがどうなったときか。また、離れる瞬間の B の速さ $v_1$ を求めよ。\pts{14}
\item 板から離れたあと、B は、離れた位置からさらに何 $\mathrm{m}$ 上昇するか。\pts{8}
\item B が板から離れないためには、押し下げる距離 $d$ をいくら以下にしなければならないか。$M,m,k,g$ で表し、数値も求めよ。\pts{10}
\item $d=0.25\unit{m}$ の場合に、B が板から受ける垂直抗力 $N$ を、位置 $x$ の関数として式で表し、B が板から離れるまでの範囲でグラフに描け。両端の値を、軸に書き入れること。\pts{13}
\end{qlist}
\end{ascolorbox4A}
% ==================================================================
%  解答・解説編　セット1
% ==================================================================
\clearpage
\begin{simple}[セット1]{解答・解説編　凡例}
各大問の冒頭の枠に、分野タグ・検算の状況・答の一覧を示す。解説は「着眼点 → 使う公式 → 解答・解説 → 確認ポイント → 誤答しやすい点」の順に述べる。配点・目標時間は問題編の各大問の見出しに記した。
\end{simple}

% ---------------------------- 解答 大問1 ----------------------------
\begin{ascolorbox4A}[75点・30分]{大問1　解答・解説}
\textgt{分野タグ}　\#熱量保存　\#理想気体　\#熱力学第一法則　\#p-Vグラフ　\#熱効率　\#分子運動論\par
\textgt{検算}　済（数式処理で全設問を再計算し、次元・単位・サイクルの熱収支を確認した）\par
\textgt{答}　(1) $5.9\times10^{3}\unit{J}$　(2) $0.84\unit{J/(g\cdot K)}$　(3) $84\unit{J/K}$　(4) $T_\mathrm{A}=\dfrac{p_0V_0}{nR}$、$T_\mathrm{B}=2T_\mathrm{A}$、$T_\mathrm{C}=4T_\mathrm{A}$、$T_\mathrm{D}=2T_\mathrm{A}$　(5) $W=p_0V_0$　(6) $Q_\mathrm{AB}=\dfrac{3}{2}p_0V_0$、$Q_\mathrm{BC}=5p_0V_0$　(7) 放出する　(8) $e=\dfrac{2}{13}$　(9) 2 倍
\end{ascolorbox4A}
\begin{multicols*}{2}
\ascboxZ{着眼点}\par
Ⅰは、高温の物体が失った熱量と低温の物体が得た熱量が等しいという熱量保存だけで解ける。Ⅱは、各過程について内部エネルギーの変化 $\Delta U$、気体がした仕事 $W$、吸収した熱量 $Q$ を第一法則で結び、状態方程式で温度を $p,V$ から決める。定積変化では $W=0$、定圧変化では $W=p\Delta V$ である。
\begin{ptbs}{KEY}[使う公式]
\begin{itemize}[leftmargin=1.2\zw,itemsep=0pt,topsep=0pt,parsep=0pt]
\item 熱量 $Q=mc\Delta T$、熱容量 $C=mc$
\item 状態方程式 $pV=nRT$、$U=\dfrac{3}{2}nRT$
\item 第一法則 $Q=\Delta U+W$（$W$ は気体がした仕事）
\item 熱効率 $e=\dfrac{W_\text{正味}}{Q_\text{吸収}}$、$v_\text{rms}\propto\sqrt{T}$
\end{itemize}
\end{ptbs}
\ascboxZ{解答・解説}\par
\textbf{(1)}\ 水が得た熱量は
\[ Q=140\times4.2\times(30-20)=5880\unit{J} \]
有効数字 2 桁で $5.9\times10^{3}\unit{J}$\ans

\textbf{(2)}\ 金属が失った熱量も $Q$ に等しいから、
\[ 100\times c\times(100-30)=5880 \]
\[ c=\frac{5880}{7000}=0.84\unit{J/(g\cdot K)}\ansm \]

\textbf{(3)}\ 熱容量は $mc=100\times0.84=84\unit{J/K}$\ans

\textbf{(4)}\ 物質量が一定であるから、状態方程式より $T\propto pV$ である。状態 A では $p_0V_0=nRT_\mathrm{A}$ より $T_\mathrm{A}=\dfrac{p_0V_0}{nR}$\ans\ 
B$(V_0,2p_0)$ では $pV$ が A の 2 倍、C$(2V_0,2p_0)$ では 4 倍、D$(2V_0,p_0)$ では 2 倍であるから
\[ T_\mathrm{B}=2T_\mathrm{A},\quad T_\mathrm{C}=4T_\mathrm{A},\quad T_\mathrm{D}=2T_\mathrm{A}\ansm \]

\textbf{(5)}\ 時計回りのサイクルであるから、正味の仕事は $p$–$V$ 図で囲まれた長方形の面積に等しい。
\[ W=(2V_0-V_0)(2p_0-p_0)=p_0V_0\ansm \]
別解として、B→C で $2p_0\times V_0$、D→A で $-p_0\times V_0$、定積変化では $0$ であるから、和は $p_0V_0$ となる。

\textbf{(6)}\ A→B は定積変化で $W=0$ である。
\[ Q_\mathrm{AB}=\Delta U=\frac{3}{2}nR(T_\mathrm{B}-T_\mathrm{A})=\frac{3}{2}p_0V_0\ansm \]
B→C は定圧変化で、気体は $W_\mathrm{BC}=2p_0\times V_0=2p_0V_0$ の仕事をする。
\[ \Delta U=\frac{3}{2}nR(T_\mathrm{C}-T_\mathrm{B})=3p_0V_0 \]
\[ Q_\mathrm{BC}=\Delta U+W_\mathrm{BC}=5p_0V_0\ansm \]

\textbf{(7)}\ D→A は定圧圧縮であるから、気体は外部から仕事をされ、$W_\mathrm{DA}=-p_0V_0<0$ である。また温度は $2T_\mathrm{A}$ から $T_\mathrm{A}$ へ下がるので $\Delta U=-\dfrac{3}{2}p_0V_0<0$ である。よって
\[ Q_\mathrm{DA}=\Delta U+W_\mathrm{DA}=-\frac{5}{2}p_0V_0<0 \]
となり、気体は熱を放出する\ans

\textbf{(8)}\ 熱を吸収するのは A→B と B→C だけである（C→D は $Q_\mathrm{CD}=-3p_0V_0$、D→A は $-\dfrac{5}{2}p_0V_0$ で放熱）。
\[ Q_\text{吸収}=\frac{3}{2}p_0V_0+5p_0V_0=\frac{13}{2}p_0V_0 \]
\[ e=\frac{W}{Q_\text{吸収}}=\frac{p_0V_0}{\frac{13}{2}p_0V_0}=\frac{2}{13}\ \ (\fallingdotseq0.15)\ansm \]

\textbf{(9)}\ 二乗平均速度は $\sqrt{\dfrac{3kT}{m}}$ に比例し、$\sqrt{T}$ に比例する。
\[ \frac{v_\mathrm{C}}{v_\mathrm{A}}=\sqrt{\frac{T_\mathrm{C}}{T_\mathrm{A}}}=\sqrt{4}=2\ \text{倍}\ansm \]
\ascboxZ{確認ポイント}\par
サイクル全体では $\Delta U=0$ であるから、$Q$ の総和 $\dfrac{3}{2}+5-3-\dfrac{5}{2}=1$ が $W=p_0V_0$ と一致する。これで (5)〜(7) の整合が確かめられる。
\begin{simple}[補足]{誤答しやすい点}
\begin{itemize}[leftmargin=1.2\zw,itemsep=0pt,topsep=0pt,parsep=0pt]
\item 定圧変化の $Q$ を $\Delta U$ だけで答えてしまう（$W=p\Delta V$ を足し忘れる）。
\item D→A を「圧縮だから熱を吸収する」と考える。温度が下がることも考えよ。
\item 熱効率の分母に、放出した熱量まで含めてしまう。
\item (1) を $5.88\times10^3$ のまま 3 桁で答えると、(2) の割り算と食い違いやすい。
\end{itemize}
\end{simple}
\end{multicols*}

% ---------------------------- 解答 大問2 ----------------------------
\begin{ascolorbox4A}[75点・30分]{大問2　解答・解説}
\textgt{分野タグ}　\#正弦波　\#波形の移動　\#気柱共鳴　\#開口端補正　\#ドップラー効果\par
\textgt{検算}　済（波形は式 $y=2.0\sin\{2\pi(x-vt)/\lambda\}$ で数値再計算し、共鳴位置は $L_1,L_2$ の両方で再確認した）\par
\textgt{答}　(1) $A=2.0\unit{cm}$、$\lambda=4.0\unit{m}$、$T=2.0\unit{s}$、$f=0.50\unit{Hz}$　(2) 正の向き　(3) 図　(4) $0.50\unit{s}$ 後　(5) $\lambda=0.500\unit{m}$、$V=340\unit{m/s}$　(6) $\Delta L=0.015\unit{m}$　(7) $L_3=0.610\unit{m}$　(8) $765\unit{Hz}$　(9) $680\unit{Hz}$　(10) 解説参照
\end{ascolorbox4A}
\begin{multicols*}{2}
\ascboxZ{着眼点}\par
Ⅰは「波形を進行方向へずらす」という見方だけで、媒質の動きも次の時刻の波形も読める。Ⅱの閉管では、閉端が節、開口端（の少し外側）が腹になる。Ⅲのドップラー効果は、音源が動くと波長が、観測者が動くと相対速さが変わる、と整理する。
\begin{ptbs}{KEY}[使う公式]
\begin{itemize}[leftmargin=1.2\zw,itemsep=0pt,topsep=0pt,parsep=0pt]
\item $v=f\lambda$、$T=\dfrac{1}{f}$
\item 閉管の共鳴 $L+\Delta L=\dfrac{(2m-1)\lambda}{4}$
\item ドップラー効果 $f'=\dfrac{V-v_\mathrm{o}}{V-v_\mathrm{s}}f_0$（$v_\mathrm{o},v_\mathrm{s}$ は音源から観測者へ向かう向きを正とした観測者・音源の速度）
\end{itemize}
\end{ptbs}
\ascboxZ{解答・解説}\par
\textbf{(1)}\ 図 2 より振幅 $A=2.0\unit{cm}$、波長 $\lambda=4.0\unit{m}$ である。
\[ T=\frac{\lambda}{v}=\frac{4.0}{2.0}=2.0\unit{s},\quad f=\frac{1}{T}=0.50\unit{Hz}\ansm \]

\textbf{(2)}\ 波は $+x$ 向きに進むので、$t=0$ の波形をわずかに $+x$ 向きへずらしたものが直後の波形である。$x=2.0\unit{m}$ のすぐ左側（$x<2.0$）では $y>0$ であるから、ずらすと $x=2.0\unit{m}$ の変位は正になる。したがって、媒質は $y$ 軸の正の向きに動く\ans

\textbf{(3)}\ $0.50\unit{s}$ の間に波は $vt=1.0\unit{m}=\dfrac{\lambda}{4}$ だけ進む。波形全体を $+x$ 向きに $1.0\unit{m}$ ずらして描けばよい。式では
\[ y=2.0\sin\frac{\pi(x-1.0)}{2}=-2.0\cos\frac{\pi x}{2} \]
となり、$x=0,1,2,3,4\unit{m}$ で $y=-2.0,0,+2.0,0,-2.0\unit{cm}$ を通る（図 8）\ans
\begin{center}
\begin{tikzpicture}
\begin{axis}[width=5.9cm,height=3.0cm,axis lines=middle,
  xlabel={$x$},ylabel={$y$},
  xlabel style={at={(ticklabel* cs:1.0)},anchor=west},
  ylabel style={at={(ticklabel* cs:1.0)},anchor=south},
  xmin=-0.3,xmax=4.8,ymin=-2.9,ymax=3.1,
  xtick={1,2,3,4},ytick={-2,2},tick label style={font=\scriptsize},label style={font=\scriptsize}]
  \addplot[densely dashed,gray,domain=0:4,samples=80] {2*sin(deg(pi*x/2))};
  \addplot[thick,domain=0:4,samples=80] {-2*cos(deg(pi*x/2))};
  \draw[->,thick] (axis cs:0.3,2.6) -- (axis cs:1.3,2.6) node[right] {\scriptsize $v$};
\end{axis}
\end{tikzpicture}\\[1pt]{\small 図8　破線：$t=0$、実線：$t=0.50\unit{s}$}
\end{center}

\textbf{(4)}\ $t=0$ に $x=0$ にある（変位 $0$ の）部分の波形が、$x=1.0\unit{m}$ へ到達したとき、その点の変位が $0$ になる。
\[ t=\frac{1.0}{2.0}=0.50\unit{s}\ansm \]
(3) の波形で $x=1.0\unit{m}$ の変位が $0$ であることとも一致する。

\textbf{(5)}\ 閉管では、隣り合う共鳴の管の長さの差が $\dfrac{\lambda}{2}$ である。
\[ \lambda=2(L_2-L_1)=2(0.360-0.110)=0.500\unit{m} \]
\[ V=f\lambda=680\times0.500=340\unit{m/s}\ansm \]

\textbf{(6)}\ 1 回目の共鳴では $L_1+\Delta L=\dfrac{\lambda}{4}$ であるから
\[ \Delta L=0.125-0.110=0.015\unit{m}\ansm \]
2 回目は $L_2+\Delta L=0.375=\dfrac{3\lambda}{4}$ で、確かに成り立つ。

\textbf{(7)}\ 3 回目の共鳴は $L_3+\Delta L=\dfrac{5\lambda}{4}=0.625\unit{m}$ より
\[ L_3=0.625-0.015=0.610\unit{m}\ansm \]
\begin{center}
\begin{tikzpicture}[line width=0.6pt,>={Stealth[length=1.6mm]}]
  \draw (-4.32,0.6) -- (0,0.6); \draw (-4.32,-0.6) -- (0,-0.6);
  \draw[very thick] (0,-0.75) -- (0,0.75);
  \draw[thick,domain=0:4.5,samples=70,smooth,variable=\u] plot ({-\u},{0.5*sin(60*\u)});
  \draw[thick,domain=0:4.5,samples=70,smooth,variable=\u] plot ({-\u},{-0.5*sin(60*\u)});
  \draw[densely dashed] (-4.5,-0.9) -- (-4.5,0.9);
  \node[above] at (0,0.75) {\scriptsize 節};
  \node[above] at (-3,0.62) {\scriptsize 節};
  \node[above] at (-1.5,0.55) {\scriptsize 腹};
  \node[above] at (-4.5,0.9) {\scriptsize 腹};
  \draw[<->] (-4.32,-1.05) -- (0,-1.05) node[midway,below] {\scriptsize $L_2$};
  \draw[<->] (-4.5,-1.05) -- (-4.32,-1.05);
  \node[left] at (-4.55,-1.05) {\scriptsize $\Delta L$};
\end{tikzpicture}\\[1pt]{\small 図9　$L=L_2$ のときの定常波}
\end{center}

\textbf{(8)}\ 音源が観測者に近づくとき、
\[ f_1=\frac{V}{V-v_\mathrm{s}}f_0=\frac{340}{340-20}\times720=765\unit{Hz}\ansm \]

\textbf{(9)}\ 音源が遠ざかるとき、
\[ f_2=\frac{V}{V+v_\mathrm{s}}f_0=\frac{340}{340+20}\times720=680\unit{Hz}\ansm \]

\textbf{(10)}\ 音源が動く場合は、音源が 1 周期の間に進む距離だけ波長が縮み、$\lambda'=\dfrac{V-v_\mathrm{s}}{f_0}$ となる。観測者は静止しているから、音は速さ $V$ で届き、$f_1=\dfrac{V}{\lambda'}$ である。一方、観測者が動く場合は、音源が静止しているので波長は $\lambda=\dfrac{V}{f_0}$ のまま変わらない。変わるのは観測者から見た音の相対速さで、$V+u$ となるから
\[ f'=\frac{V+u}{\lambda}=\frac{340+20}{340}\times720\fallingdotseq762\unit{Hz} \]
となる。このように、波長が変わるか、相対速さが変わるかの違いにより、$f_1=765\unit{Hz}$ とはわずかに異なる\ans
\ascboxZ{確認ポイント}\par
(5)〜(7) は、$L_1$ と $L_2$ のどちらを使っても同じ $\Delta L$ になるかで検算できる。(8)(9) は、近づくと高く、遠ざかると低くなる向きになっている。
\begin{simple}[補足]{誤答しやすい点}
\begin{itemize}[leftmargin=1.2\zw,itemsep=0pt,topsep=0pt,parsep=0pt]
\item (2) で、波の進む向きと媒質の動く向きを混同する（媒質は波の進行方向とは別に上下へ動く）。
\item (5) で $L_1$ を $\dfrac{\lambda}{4}$ とみなして $\lambda=4L_1$ とする（開口端補正を無視した誤り）。
\item (8)(9) で、分母の符号を逆にする。近づく音源は波長を縮める。
\end{itemize}
\end{simple}
\end{multicols*}

% ---------------------------- 解答 大問3 ----------------------------
\begin{ascolorbox4A}[75点・30分]{大問3　解答・解説}
\textgt{分野タグ}　\#直流回路　\#合成抵抗　\#ジュール熱　\#コンデンサー　\#キルヒホッフの法則　\#過渡現象（概形）\par
\textgt{検算}　済（回路方程式を数値で解き直し、$S_2$ を開く前後の電流の和 $I_1'=I_3+I_C$ を確認した）\par
\textgt{答}　(1) $8.0\,\Omega$、$I=3.0\unit{A}$　(2) $I_2=2.0\unit{A}$、$I_3=1.0\unit{A}$　(3) $P_1=36\unit{W}$、$H=360\unit{J}$　(4) $V_C=12\unit{V}$、$Q=2.4\times10^{-5}\unit{C}$、$U=1.4\times10^{-4}\unit{J}$　(5) $I_1'=3.0\unit{A}$、$I_C=2.0\unit{A}$　(6) $V_C'=18\unit{V}$、$Q'=3.6\times10^{-5}\unit{C}$、$\Delta Q=1.2\times10^{-5}\unit{C}$　(7) 図 10
\end{ascolorbox4A}
\begin{multicols*}{2}
\ascboxZ{着眼点}\par
コンデンサーは、定常状態では電流を通さず（断線と同じ）、電荷が変化する過程では、その瞬間の電圧を保ったまま電流が流れ込む。そこで、「スイッチを操作した直後」と「十分に時間がたった後」を分けて、それぞれの電流の分布を考える。
\begin{ptbs}{KEY}[使う公式]
\begin{itemize}[leftmargin=1.2\zw,itemsep=0pt,topsep=0pt,parsep=0pt]
\item オームの法則 $V=RI$、並列 $\dfrac{1}{R}=\dfrac{1}{R_2}+\dfrac{1}{R_3}$
\item 電力 $P=IV=I^2R$、ジュール熱 $H=Pt$
\item コンデンサー $Q=CV$、$U=\dfrac{1}{2}CV^2$
\item キルヒホッフの第 1・第 2 法則
\end{itemize}
\end{ptbs}
\ascboxZ{解答・解説}\par
\textbf{(1)}\ $R_2$ と $R_3$ の並列部分の合成抵抗 $R_\mathrm{p}$ は
\[ \frac{1}{R_\mathrm{p}}=\frac{1}{6.0}+\frac{1}{12}=\frac{1}{4.0}\ \ \therefore\ R_\mathrm{p}=4.0\,\Omega \]
全体の合成抵抗は $R_1+R_\mathrm{p}=4.0+4.0=8.0\,\Omega$ である。
\[ I=\frac{E}{8.0}=\frac{24}{8.0}=3.0\unit{A}\ansm \]

\textbf{(2)}\ 並列部分の電圧は $V_\mathrm{AB}=IR_\mathrm{p}=3.0\times4.0=12\unit{V}$ である。
\[ I_2=\frac{12}{6.0}=2.0\unit{A},\quad I_3=\frac{12}{12}=1.0\unit{A}\ansm \]
和は $3.0\unit{A}$ で、(1) の $I$ に一致する。

\textbf{(3)}\ 
\[ P_1=I^2R_1=3.0^2\times4.0=36\unit{W} \]
\[ H=P_1t=36\times10=360\unit{J}\ansm \]

\textbf{(4)}\ 十分に時間がたつとコンデンサーに電流は流れないから、$\mathrm{S}_1$ を閉じても (1)(2) の電流分布は変わらない。コンデンサーは A–B 間に並列につながるので
\[ V_C=V_\mathrm{AB}=12\unit{V} \]
\[ Q=CV_C=2.0\times10^{-6}\times12=2.4\times10^{-5}\unit{C} \]
\[ U=\frac{1}{2}CV_C^2=\frac{1}{2}\times2.0\times10^{-6}\times12^2=1.44\times10^{-4}\unit{J} \]
有効数字 2 桁で $1.4\times10^{-4}\unit{J}$\ans

\textbf{(5)}\ $\mathrm{S}_2$ を開いた直後も、コンデンサーの電気量は変わらないから $V_C=12\unit{V}$ のままである。これは $V_\mathrm{AB}$ に等しい。
\[ I_3=\frac{12}{12}=1.0\unit{A} \]
$R_1$ の両端の電圧は $E-V_\mathrm{AB}=24-12=12\unit{V}$ であるから
\[ I_1'=\frac{12}{4.0}=3.0\unit{A}\ansm \]
点 A で電流の法則を用いると $I_1'=I_3+I_C$ であるから
\[ I_C=3.0-1.0=2.0\unit{A}\ansm \]

\textbf{(6)}\ 十分に時間がたつと $I_C=0$ となり、電流は $R_1$ と $R_3$ を直列に流れる。
\[ I=\frac{24}{4.0+12}=1.5\unit{A} \]
\[ V_C'=IR_3=1.5\times12=18\unit{V}\ansm \]
\[ Q'=CV_C'=2.0\times10^{-6}\times18=3.6\times10^{-5}\unit{C}\ansm \]
\[ \Delta Q=Q'-Q=1.2\times10^{-5}\unit{C}\ansm \]

\textbf{(7)}\ $V_C$ は $t=0$ で $12\unit{V}$ であり、そこから増加して $18\unit{V}$ に近づく。$t=0$ の傾きは、$I_C=C\dfrac{dV_C}{dt}$ から最も大きく、時間とともに充電電流が減るので傾きは小さくなる（上に凸の曲線）。$18\unit{V}$ を超えることはない（図 10）\ans
\begin{center}
\begin{tikzpicture}[>={Stealth[length=1.8mm]},line width=0.6pt,x=1.0cm,y=0.24cm]
  \draw[->] (0,0) -- (5.6,0) node[right] {\small $t$};
  \draw[->] (0,0) -- (0,21) node[above] {\small $V_C\,[\mathrm{V}]$};
  \draw[densely dashed] (0,18) node[left] {\small $18$} -- (5.4,18);
  \draw[densely dashed] (0,12) node[left] {\small $12$} -- (0.3,12);
  \draw[thick,domain=0:5.2,samples=60,smooth] plot ({\x},{18-6*exp(-\x)});
  \fill (0,12) circle[radius=0.05];
  \node[below left] at (0,0) {\small O};
\end{tikzpicture}\\[1pt]{\small 図10}
\end{center}
\ascboxZ{確認ポイント}\par
$S_2$ を開く前の $V_C=12\unit{V}$、開いた後の $V_C'=18\unit{V}$ は、いずれも「コンデンサーの両端の電圧 $=$ 並列につながる抵抗の電圧」として同じ方針で求められる。$V_C'$ が $E=24\unit{V}$ を超えないことも確かめておく。
\begin{simple}[発展]{充電の速さ（参考）}
充電中のコンデンサーから見ると、電池と $R_1$ を、$R_1$ と $R_3$ の並列 $R=\dfrac{4.0\times12}{4.0+12}=3.0\,\Omega$ として扱える。このとき $V_C=18-6e^{-t/\tau}$（$\tau=RC=6.0\,\mu\mathrm{s}$）となる。$t=0$ の傾き $6/\tau=1.0\times10^{6}\unit{V/s}$ は、(5) の $I_C/C=\dfrac{2.0}{2.0\times10^{-6}}$ と一致する。
\end{simple}
\begin{simple}[補足]{誤答しやすい点}
\begin{itemize}[leftmargin=1.2\zw,itemsep=0pt,topsep=0pt,parsep=0pt]
\item $V_C$ を電池の起電力 $24\unit{V}$ としてしまう。コンデンサーにかかるのは $R_1$ の電圧降下を引いた $V_\mathrm{AB}$ である。
\item $\mathrm{S}_2$ を開いた直後に $V_C$ が瞬時に $18\unit{V}$ になると考える。電荷は連続的にしか変化しない。
\item 静電エネルギーの $\dfrac{1}{2}$ を落とす、$\mu\mathrm{F}$ を $10^{-6}\unit{F}$ に直し忘れる。
\end{itemize}
\end{simple}
\end{multicols*}

% ---------------------------- 解答 大問4 ----------------------------
\begin{ascolorbox4A}[75点・30分]{大問4　解答・解説}
\textgt{分野タグ}　\#エネルギー保存　\#単振動　\#運動量保存　\#完全非弾性衝突　\#鉛直円運動　\#放物運動\par
\textgt{検算}　済（全設問を記号計算で導出し直し、数値例 $m=1.0\unit{kg}$、$k=400\unit{N/m}$、$d=0.50\unit{m}$、$R=0.20\unit{m}$ でエネルギー保存と $N_\mathrm{Q}$ の符号を確認した）\par
\textgt{答}　(1) $v_0=d\sqrt{k/m}$　(2) $t_1=\dfrac{\pi}{2}\sqrt{m/k}$　(3) $v_1=\dfrac{d}{2}\sqrt{k/m}$、$\Delta E=\dfrac{kd^2}{4}$　(4) $N_\mathrm{P}=2mg+\dfrac{kd^2}{2R}$　(5) $v_\mathrm{Q}=\sqrt{v_1^2-4gR}$、$N_\mathrm{Q}=\dfrac{2mv_1^2}{R}-10mg$　(6) $d\geqq 2\sqrt{\dfrac{5mgR}{k}}$　(7) $t=2\sqrt{\dfrac{R}{g}}$、$x=2v_\mathrm{Q}\sqrt{\dfrac{R}{g}}$　(8) $d=2\sqrt{\dfrac{7mgR}{k}}$
\end{ascolorbox4A}
\begin{multicols*}{2}
\ascboxZ{着眼点}\par
物体の運動を「ばねによる加速 → 衝突 → 円運動 → 放物運動」と 4 つの段階に分け、それぞれの段階で成り立つ保存則を選ぶ。ばねの段階は力学的エネルギー保存と単振動、衝突は運動量保存（エネルギーは保存しない）、なめらかな半円筒の上りは力学的エネルギー保存と円運動の運動方程式、飛び出した後は放物運動である。
\begin{ptbs}{KEY}[使う公式]
\begin{itemize}[leftmargin=1.2\zw,itemsep=0pt,topsep=0pt,parsep=0pt]
\item 弾性エネルギー $\dfrac{1}{2}kx^2$、周期 $T=2\pi\sqrt{\dfrac{m}{k}}$
\item 運動量保存 $mv_0=(m+m)v_1$
\item 円運動 $m\dfrac{v^2}{R}=$（向心方向の力の和）
\item 面から離れない条件 $N\geqq0$
\end{itemize}
\end{ptbs}
\ascboxZ{解答・解説}\par
\textbf{(1)}\ なめらかな床であるから、ばねの弾性エネルギーがすべて A の運動エネルギーになる。
\[ \frac{1}{2}kd^2=\frac{1}{2}mv_0^2\ \ \therefore\ v_0=d\sqrt{\frac{k}{m}}\ansm \]

\textbf{(2)}\ ばねについた A は、自然の長さの位置 $\mathrm{O}$ を中心とする、振幅 $d$、周期 $2\pi\sqrt{\dfrac{m}{k}}$ の単振動をする。最大の縮みの位置から中心 $\mathrm{O}$ までは $\dfrac{1}{4}$ 周期である。$\mathrm{O}$ でばねの力が $0$ になり、以後 A は等速でばねから離れていく。
\[ t_1=\frac{1}{4}\times2\pi\sqrt{\frac{m}{k}}=\frac{\pi}{2}\sqrt{\frac{m}{k}}\ansm \]

\textbf{(3)}\ 衝突では A、B の系の運動量が保存される。
\[ mv_0=(m+m)v_1\ \ \therefore\ v_1=\frac{v_0}{2}=\frac{d}{2}\sqrt{\frac{k}{m}}\ansm \]
失われた力学的エネルギーは
\[ \Delta E=\frac{1}{2}mv_0^2-\frac{1}{2}(2m)v_1^2=\frac{1}{4}mv_0^2=\frac{kd^2}{4}\ansm \]
ばねの弾性エネルギー $\dfrac{1}{2}kd^2$ の半分が失われている。

\textbf{(4)}\ P は円の最下点であるから、向心方向は鉛直上向きである。質量 $2m$ の運動方程式は
\[ 2m\frac{v_1^2}{R}=N_\mathrm{P}-2mg \]
$v_1^2=\dfrac{kd^2}{4m}$ を代入して
\[ N_\mathrm{P}=2mg+\frac{2mv_1^2}{R}=2mg+\frac{kd^2}{2R}\ansm \]

\textbf{(5)}\ P から Q へ、なめらかな面の上で、高さが $2R$ だけ上がる。力学的エネルギー保存より
\[ \frac{1}{2}(2m)v_1^2=\frac{1}{2}(2m)v_\mathrm{Q}^2+(2m)g(2R) \]
\[ v_\mathrm{Q}=\sqrt{v_1^2-4gR}\ansm \]
Q では向心方向が鉛直下向きで、重力 $2mg$ と垂直抗力 $N_\mathrm{Q}$ がともに向心方向にはたらく。
\[ 2m\frac{v_\mathrm{Q}^2}{R}=N_\mathrm{Q}+2mg \]
\[ N_\mathrm{Q}=\frac{2m(v_1^2-4gR)}{R}-2mg=\frac{2mv_1^2}{R}-10mg\ansm \]

\textbf{(6)}\ 半円筒の内側面から離れないためには、Q で $N_\mathrm{Q}\geqq0$ であればよい（途中の点では垂直抗力はこれより大きい）。
\[ \frac{2mv_1^2}{R}\geqq10mg\ \ \therefore\ v_1^2\geqq5gR \]
$v_1^2=\dfrac{kd^2}{4m}$ を代入して
\[ d^2\geqq\frac{20mgR}{k}\ \ \therefore\ d\geqq2\sqrt{\frac{5mgR}{k}}\ansm \]

\textbf{(7)}\ Q から水平に飛び出し、高さ $2R$ を落下する。
\[ 2R=\frac{1}{2}gt^2\ \ \therefore\ t=2\sqrt{\frac{R}{g}}\ansm \]
水平方向は等速であるから
\[ x=v_\mathrm{Q}t=2v_\mathrm{Q}\sqrt{\frac{R}{g}}\ansm \]

\textbf{(8)}\ $x=2\sqrt{3}R$ を (7) に代入すると
\[ v_\mathrm{Q}=\frac{x}{t}=\frac{2\sqrt{3}R}{2\sqrt{R/g}}=\sqrt{3gR},\quad v_\mathrm{Q}^2=3gR \]
(5) より $v_1^2=v_\mathrm{Q}^2+4gR=7gR$ である。$v_0=2v_1$ であるから $v_0^2=28gR$ となり、(1) より
\[ d^2=\frac{mv_0^2}{k}=\frac{28mgR}{k}\ \ \therefore\ d=2\sqrt{\frac{7mgR}{k}}\ansm \]
\ascboxZ{確認ポイント}\par
(8) の $d$ は (6) の条件 $d\geqq2\sqrt{5mgR/k}$ を満たす（$7>5$）。このとき $N_\mathrm{Q}=\dfrac{2m\cdot7gR}{R}-10mg=4mg>0$ で、Q で面から離れない。
\begin{simple}[補足]{誤答しやすい点}
\begin{itemize}[leftmargin=1.2\zw,itemsep=0pt,topsep=0pt,parsep=0pt]
\item (3) の衝突で、力学的エネルギー保存を用いてしまう（完全非弾性衝突では保存しない）。
\item 一体の物体の質量は $2m$ である。$N_\mathrm{P}$、$N_\mathrm{Q}$、位置エネルギーで $m$ のままにしがち。
\item Q での重力は向心方向にはたらく。$N_\mathrm{Q}-2mg$ と符号を誤らない。
\item (6) を「Q で速さ $0$」と考える。面の上では、離れない条件は $N_\mathrm{Q}\geqq0$ である。
\end{itemize}
\end{simple}
\end{multicols*}
% ==================================================================
%  解答・解説編　セット2
% ==================================================================
\clearpage
\begin{simple}[セット2]{解答・解説編　凡例}
各大問の冒頭の枠に、分野タグ・検算の状況・答の一覧を示す。解説は「着眼点 → 使う公式 → 解答・解説 → 確認ポイント → 誤答しやすい点」の順に述べる。
\end{simple}

% ---------------------------- 解答 大問1 ----------------------------
\begin{ascolorbox4A}[75点・30分]{大問1　解答・解説}
\textgt{分野タグ}　\#斜面　\#動摩擦力　\#静止摩擦力　\#運動方程式　\#等速円運動　\#円錐振り子\par
\textgt{検算}　済（斜面は数値で往復の運動を再計算し、仕事とエネルギーの関係で $v$ を別解と照合した。円錐振り子は $\theta=60^\circ$、$L=4.9\unit{m}$ で $T=\pi\unit{s}$ を数値確認した）\par
\textgt{答}　(1) 図 11　(2) $a_1=\dfrac{3}{4}g$　(3) $s=\dfrac{2v_0^2}{3g}$、$3.3\unit{m}$　(4) 滑り降りる　(5) $v=\dfrac{v_0}{\sqrt{3}}$　(6) $S=\dfrac{mg}{\cos\theta}$、$v=\sin\theta\sqrt{\dfrac{gL}{\cos\theta}}$　(7) $T=2\pi\sqrt{\dfrac{L\cos\theta}{g}}$　(8) $T=\pi\unit{s}\fallingdotseq3.1\unit{s}$、$S=39\unit{N}$　(9) 短くなる
\end{ascolorbox4A}
\begin{multicols*}{2}
\ascboxZ{着眼点}\par
Ⅰは、斜面に平行な方向と垂直な方向に力を分解する。上るときと下るときで、動摩擦力の向きが逆になる点に注意する。Ⅱは、鉛直方向の力のつり合いと、水平方向（円の中心向き）の運動方程式を立てる。
\begin{ptbs}{KEY}[使う公式]
\begin{itemize}[leftmargin=1.2\zw,itemsep=0pt,topsep=0pt,parsep=0pt]
\item 動摩擦力 $f'=\mu'N$、最大静止摩擦力 $f_0=\mu N$
\item 等加速度運動 $v^2-v_0^2=2ax$
\item 等速円運動 $m\dfrac{v^2}{r}=$（向心方向の力）、$T=\dfrac{2\pi r}{v}$
\end{itemize}
\end{ptbs}
\ascboxZ{解答・解説}\par
\textbf{(1)}\ 重力 $mg$（鉛直下向き）、垂直抗力 $N$（斜面に垂直で上向き）、動摩擦力 $\mu'N$（物体は斜面を上向きに滑るので、斜面に沿って下向き）の 3 力がはたらく（図 11）。摩擦力は運動の向きと逆であり、打ち出した力は物体にはたらき続けないことに注意する\ans
\begin{center}
\begin{tikzpicture}[>={Stealth[length=2mm]},line width=0.6pt,force/.style={->,very thick},x=0.72cm,y=0.72cm]
  \def\ang{30}
  \fill[pattern={Lines[angle=45,distance=2mm]}] (0,0) -- (6,0) -- (6,-0.25) -- (0,-0.25) -- cycle;
  \draw (0,0) -- (6,0) -- (6,{6*tan(\ang)}) -- cycle;
  \begin{scope}[rotate=\ang]
    \draw[fill=gray!20] (1.6,0) rectangle (2.7,0.8);
    \coordinate (G) at (2.15,0.4);
    \draw[force] (G) -- ++(0,1.4) node[above left] {\small $N$};
    \draw[force] (G) -- ++(-1.4,0) node[left] {\small $\mu'N$};
  \end{scope}
  \draw[force] (G) -- ++(0,-1.5) node[below] {\small $mg$};
\end{tikzpicture}\\[1pt]{\small 図11}
\end{center}

\textbf{(2)}\ 斜面に垂直な方向のつり合いより $N=mg\cos\theta$ である。斜面に沿って上向きを正とすると、重力の成分 $mg\sin\theta$ と動摩擦力はともに下向きである。
\[ ma=-mg\sin\theta-\mu'mg\cos\theta \]
\[ a_1=g\left(\frac{1}{2}+\frac{\sqrt{3}}{6}\cdot\frac{\sqrt{3}}{2}\right)=\frac{3}{4}g\ansm \]

\textbf{(3)}\ 最高点で速さが $0$ になるから、$0-v_0^2=-2a_1s$ より
\[ s=\frac{v_0^2}{2a_1}=\frac{2v_0^2}{3g}\ansm \]
$v_0=7.0\unit{m/s}$、$g=9.8\unit{m/s^2}$ を代入して
\[ s=\frac{2\times49}{3\times9.8}=\frac{10}{3}\fallingdotseq3.3\unit{m}\ansm \]

\textbf{(4)}\ 最大静止摩擦力は
\[ \mu N=\frac{1}{2}\cdot\frac{\sqrt{3}}{2}mg=\frac{\sqrt{3}}{4}mg\fallingdotseq0.43mg \]
一方、重力の斜面方向の成分は $mg\sin\theta=\dfrac{1}{2}mg=0.50mg$ である。重力の成分が最大静止摩擦力より大きいので、物体はそこに静止できず、斜面を滑り降りる\ans
（$\mu<\tan30^\circ=\dfrac{1}{\sqrt{3}}\fallingdotseq0.58$ と同値である。）

\textbf{(5)}\ 滑り降りるとき、動摩擦力は上向きにはたらく。
\[ a_2=g\left(\sin\theta-\mu'\cos\theta\right)=g\left(\frac{1}{2}-\frac{1}{4}\right)=\frac{g}{4} \]
距離は同じ $s$ であるから
\[ v^2=2a_2s=2\cdot\frac{g}{4}\cdot\frac{2v_0^2}{3g}=\frac{v_0^2}{3}\ \ \therefore\ v=\frac{v_0}{\sqrt{3}}\ansm \]
\begin{simple}[別解]{仕事とエネルギー}
往復の間に、動摩擦力は $2s$ の道のりで負の仕事をする。$\mu'\cos\theta=\dfrac{1}{4}$ であるから
\[ \frac{1}{2}mv^2=\frac{1}{2}mv_0^2-\frac{1}{4}mg\cdot2s \]
$s=\dfrac{2v_0^2}{3g}$ を代入すると $v^2=\dfrac{v_0^2}{3}$ となり、(5) と一致する。
\end{simple}

\textbf{(6)}\ 糸の張力 $S$ の鉛直成分が重力とつり合う。
\[ S\cos\theta=mg\ \ \therefore\ S=\frac{mg}{\cos\theta}\ansm \]
水平成分が向心力になる。円の半径は $r=L\sin\theta$ である。
\[ m\frac{v^2}{L\sin\theta}=S\sin\theta=mg\tan\theta \]
\[ v^2=gL\sin\theta\tan\theta\ \ \therefore\ v=\sin\theta\sqrt{\frac{gL}{\cos\theta}}\ansm \]

\textbf{(7)}\ 
\[ T=\frac{2\pi r}{v}=\frac{2\pi L\sin\theta}{\sin\theta\sqrt{gL/\cos\theta}} \]
\[ =2\pi\sqrt{\frac{L\cos\theta}{g}}\ansm \]

\textbf{(8)}\ $\cos60^\circ=\dfrac{1}{2}$ であるから
\[ T=2\pi\sqrt{\frac{4.9\times0.5}{9.8}}=2\pi\times\frac{1}{2}=\pi\fallingdotseq3.1\unit{s}\ansm \]
\[ S=\frac{2.0\times9.8}{0.5}=39.2\fallingdotseq39\unit{N}\ansm \]

\textbf{(9)}\ $T=2\pi\sqrt{L\cos\theta/g}$ で、$L$ が一定のまま $\theta$ が大きくなると $\cos\theta$ は小さくなる。したがって周期 $T$ は短くなる\ans
\ascboxZ{確認ポイント}\par
(7) で $\theta\to0$ とすると $T\to2\pi\sqrt{L/g}$（単振り子の周期）となり、極限の確認になる。(8) の数値は、(6) の $v$ から $T=2\pi r/v$ で求めても $\pi\unit{s}$ になる。
\begin{simple}[補足]{誤答しやすい点}
\begin{itemize}[leftmargin=1.2\zw,itemsep=0pt,topsep=0pt,parsep=0pt]
\item 上るときも下るときも同じ向きに動摩擦力を描く。摩擦力は運動の向きと逆である。
\item (4) で動摩擦係数 $\mu'$ を使ってしまう。静止の判断には静止摩擦係数 $\mu$ を使う。
\item 円錐振り子で、遠心力を持ち出しながら向心力も加えてしまう。慣性系では向心力だけを考える。
\item 円運動の半径に $L$ を使う（正しくは $r=L\sin\theta$）。
\end{itemize}
\end{simple}
\end{multicols*}

% ---------------------------- 解答 大問2 ----------------------------
\begin{ascolorbox4A}[75点・30分]{大問2　解答・解説}
\textgt{分野タグ}　\#ローレンツ力　\#荷電粒子の円運動　\#電磁誘導　\#レンツの法則　\#エネルギー保存\par
\textgt{検算}　済（$r$、$T$ は記号計算、誘導の数値は $V_0=B\ell v=1.0\unit{V}$、$P=Fv=I^2R=0.50\unit{W}$ で整合を確認した。電流の向きはレンツの法則とローレンツ力の両方で確認した）\par
\textgt{答}　(1) $v=\sqrt{\dfrac{2qV}{m}}$　(2) $r=\dfrac{1}{B}\sqrt{\dfrac{2mV}{q}}$、$T=\dfrac{2\pi m}{qB}$　(3) $2r$　(4) $\sqrt{2}$ 倍　(5) 解説参照　(6) $1.0\unit{V}$　(7) $0.50\unit{A}$、P→Q、Q　(8) $F=0.10\unit{N}$、$P=0.50\unit{W}$　(9) $P_R=0.50\unit{W}$（$=P$）　(10) 減速して $0$ に近づく
\end{ascolorbox4A}
\begin{multicols*}{2}
\ascboxZ{着眼点}\par
Ⅰは、加速（電場がする仕事）と、磁場中の円運動（ローレンツ力が向心力）の 2 段階に分ける。Ⅱは、回路を貫く磁束の変化から起電力を求め、電流の向きをレンツの法則で決める。棒が磁場から受ける力が、外力の仕事率とジュール熱をつなぐ。
\begin{ptbs}{KEY}[使う公式]
\begin{itemize}[leftmargin=1.2\zw,itemsep=0pt,topsep=0pt,parsep=0pt]
\item ローレンツ力 $f=qvB$（$\vec{v}\perp\vec{B}$）
\item 加速 $qV=\dfrac{1}{2}mv^2$
\item 誘導起電力 $V_0=-\dfrac{\Delta\Phi}{\Delta t}$、棒の起電力 $V_0=B\ell v$
\item 電流が磁場から受ける力 $F=IB\ell$
\end{itemize}
\end{ptbs}
\ascboxZ{解答・解説}\par
\textbf{(1)}\ 電場が粒子にする仕事 $qV$ が運動エネルギーになる。
\[ qV=\frac{1}{2}mv^2\ \ \therefore\ v=\sqrt{\frac{2qV}{m}}\ansm \]

\textbf{(2)}\ ローレンツ力 $qvB$ が円運動の向心力になる。
\[ m\frac{v^2}{r}=qvB\ \ \therefore\ r=\frac{mv}{qB}=\frac{1}{B}\sqrt{\frac{2mV}{q}}\ansm \]
\[ T=\frac{2\pi r}{v}=\frac{2\pi m}{qB}\ansm \]
周期は速さによらない。

\textbf{(3)}\ 境界に垂直に入射した粒子は半円を描き、再び境界に垂直に出る。したがって $\mathrm{OP}=2r$\ans

\textbf{(4)}\ $r=\dfrac{1}{B}\sqrt{\dfrac{2mV}{q}}$ は $\sqrt{m}$ に比例する。質量が $2$ 倍になると、半径は $\sqrt{2}$ 倍になる\ans

\textbf{(5)}\ ローレンツ力は、常に速度に垂直な向きにはたらくから、粒子に仕事をしない。したがって運動エネルギーは変わらず、速さが一定に保たれる\ans

\textbf{(6)}\ 時間 $\Delta t$ の間に棒は $v\Delta t$ 進み、回路の面積は $\ell v\Delta t$ だけ増える。貫く磁束の変化は $\Delta\Phi=B\ell v\Delta t$ であるから
\[ V_0=\frac{\Delta\Phi}{\Delta t}=B\ell v=0.50\times0.40\times5.0=1.0\unit{V}\ansm \]

\textbf{(7)}\ 
\[ I=\frac{V_0}{R}=\frac{1.0}{2.0}=0.50\unit{A}\ansm \]
上向きに貫く磁束が増加するので、レンツの法則により、誘導電流は下向きの磁束をつくる向き、すなわち上から見て時計回りに流れる。したがって、棒の中では P→Q の向きに流れる\ans
（確認：棒の中の正の電荷には、$q\vec{v}\times\vec{B}$ より下向き、すなわち Q へ向かう力がはたらく。）棒は電源の役割をしており、電源の内部では電流は低電位側から高電位側へ流れるので、電位が高いのは Q である\ans

\textbf{(8)}\ 棒を流れる電流が磁場から受ける力は
\[ F_\mathrm{m}=IB\ell=0.50\times0.50\times0.40=0.10\unit{N} \]
で、向きは運動と逆（左向き）である。等速で動かすには、これと同じ大きさの外力が右向きに必要である。
\[ F=0.10\unit{N}\ansm \]
\[ P=Fv=0.10\times5.0=0.50\unit{W}\ansm \]

\textbf{(9)}\ 
\[ P_R=I^2R=0.50^2\times2.0=0.50\unit{W}\ansm \]
これは $P$ に等しい。外力が棒にする仕事が、すべて抵抗でジュール熱に変わっている（エネルギー保存）\ans

\textbf{(10)}\ 外力がなくなると、棒には磁場からの力 $F_\mathrm{m}=\dfrac{B^2\ell^2v}{R}$ だけが運動と逆向きにはたらくので、棒は減速する。$v$ が小さくなると $F_\mathrm{m}$ も小さくなるから、減速のしかたは次第に緩やかになり、速さは $0$ に近づく。失われた運動エネルギーは、ジュール熱になる\ans
\ascboxZ{確認ポイント}\par
(8) と (9) で、$Fv=I^2R$ が成り立つことを必ず確かめる。起電力の電力 $V_0I=1.0\times0.50=0.50\unit{W}$ とも一致する。
\begin{simple}[補足]{誤答しやすい点}
\begin{itemize}[leftmargin=1.2\zw,itemsep=0pt,topsep=0pt,parsep=0pt]
\item 電流の向きと、電位が高い側を取り違える。棒の中では電流は電位が低い方から高い方へ流れる。
\item (2) の $T$ を、$v$ や $V$ を含む式のまま答える。周期は $q,m,B$ だけで決まる。
\item (10) で、棒が等速のまま動き続けると考える。摩擦がなくても、電流が磁場から受ける力が運動を妨げる。
\end{itemize}
\end{simple}
\end{multicols*}

% ---------------------------- 解答 大問3 ----------------------------
\begin{ascolorbox4A}[75点・30分]{大問3　解答・解説}
\textgt{分野タグ}　\#ヤングの実験　\#光路差　\#屈折率と光路長　\#回折格子　\#分散\par
\textgt{検算}　済（$\Delta y$、$y_0$、$\sin\theta_m=m/4$ を数値再計算し、$(n-1)t=5\lambda$ で 5 本分のずれになることを確認した）\par
\textgt{答}　(1) 波長・速さは $\dfrac{1}{n}$ 倍、振動数は $1$ 倍　(2) $\dfrac{dy}{L}$　(3) $y_m=\dfrac{mL\lambda}{d}$、$\Delta y=\dfrac{L\lambda}{d}=2.5\times10^{-3}\unit{m}$　(4) $\dfrac{dy}{L}-(n-1)t=m\lambda$　(5) $y_0=\dfrac{L(n-1)t}{d}=1.25\times10^{-2}\unit{m}$、5 倍分、$\mathrm{S}_1$ 側　(6) $d'=2.0\times10^{-6}\unit{m}$、$\theta_2=30^\circ$　(7) 7 本　(8) 解説参照
\end{ascolorbox4A}
\begin{multicols*}{2}
\ascboxZ{着眼点}\par
干渉の問題では、まず「2 つの光が届くまでの経路差（光路差）」を式にし、それが波長の整数倍なら明線、半整数倍なら暗線という条件に結びつける。媒質中では波長が変わるので、屈折率のある物質を通る光は「光路長」で比べる。
\begin{ptbs}{KEY}[使う公式]
\begin{itemize}[leftmargin=1.2\zw,itemsep=0pt,topsep=0pt,parsep=0pt]
\item 屈折率 $n=\dfrac{c}{v}=\dfrac{\lambda_0}{\lambda}$、光路長 $=n\times$ 距離
\item ヤング：明線 $\dfrac{dy}{L}=m\lambda$
\item 回折格子：明線 $d'\sin\theta_m=m\lambda$
\end{itemize}
\end{ptbs}
\ascboxZ{解答・解説}\par
\textbf{(1)}\ 光の振動数は媒質によらず変わらない。速さは $v=\dfrac{c}{n}$ で $\dfrac{1}{n}$ 倍になるから、$v=f\lambda$ より波長も $\dfrac{1}{n}$ 倍になる\ans

\textbf{(2)}\ 2 つのスリットを $\mathrm{S}_1\left(0,\tfrac{d}{2}\right)$、$\mathrm{S}_2\left(0,-\tfrac{d}{2}\right)$、点 P を $(L,y)$ とすると
\[ \mathrm{S_2P}^2-\mathrm{S_1P}^2=\left(y+\frac{d}{2}\right)^2-\left(y-\frac{d}{2}\right)^2=2yd \]
$\mathrm{S_2P}+\mathrm{S_1P}\fallingdotseq2L$ であるから
\[ \mathrm{S_2P}-\mathrm{S_1P}=\frac{2yd}{\mathrm{S_2P}+\mathrm{S_1P}}\fallingdotseq\frac{dy}{L}\ansm \]

\textbf{(3)}\ 明線の条件は経路差が波長の整数倍になることである。
\[ \frac{dy}{L}=m\lambda\ \ \therefore\ y_m=\frac{mL\lambda}{d}\ansm \]
\[ \Delta y=\frac{L\lambda}{d} \]
\[ =\frac{1.0\times5.0\times10^{-7}}{2.0\times10^{-4}}=2.5\times10^{-3}\unit{m}\ansm \]

\textbf{(4)}\ $\mathrm{S}_1$ 側の光は、ガラス板により光路長が $(n-1)t$ だけ長くなる。したがって実質の光路差は
\[ \mathrm{S_2P}-\left\{\mathrm{S_1P}+(n-1)t\right\}=\frac{dy}{L}-(n-1)t \]
これが波長の整数倍のとき明線となる。
\[ \frac{dy}{L}-(n-1)t=m\lambda\ansm \]

\textbf{(5)}\ $m=0$ とおくと
\[ y_0=\frac{L(n-1)t}{d} \]
\[ =\frac{1.0\times0.5\times5.0\times10^{-6}}{2.0\times10^{-4}}=1.25\times10^{-2}\unit{m}\ansm \]
$y_0>0$ であるから、中央の明線は $\mathrm{S}_1$ 側（ガラス板を置いた側）へずれる。また $(n-1)t=2.5\times10^{-6}\unit{m}=5\lambda$ であるから
\[ \frac{y_0}{\Delta y}=\frac{1.25\times10^{-2}}{2.5\times10^{-3}}=5\ \text{倍}\ansm \]
明線の間隔の 5 本分だけ、$\mathrm{S}_1$ 側へずれる。

\textbf{(6)}\ 格子定数は溝の間隔であるから
\[ d'=\frac{1\times10^{-3}}{500}=2.0\times10^{-6}\unit{m}\ansm \]
$d'\sin\theta_2=2\lambda$ より
\[ \sin\theta_2=\frac{2\times5.0\times10^{-7}}{2.0\times10^{-6}}=\frac{1}{2}\ \ \therefore\ \theta_2=30^\circ\ansm \]

\textbf{(7)}\ $\sin\theta_m=\dfrac{m\lambda}{d'}=\dfrac{m}{4}$ で、$|\sin\theta_m|<1$ でなければならない。$m=\pm4$ では $\sin\theta=1$（$\theta=90^\circ$、格子面に平行）となり、スクリーン上には現れない。したがって $m=0,\pm1,\pm2,\pm3$ の 7 本である\ans

\textbf{(8)}\ 1 次の明線は $d'\sin\theta_1=\lambda$ より $\sin\theta_1=\dfrac{\lambda}{d'}$ を満たす。格子定数 $d'$ は同じであるから、波長が長い光ほど回折角 $\theta_1$ が大きい。赤色（波長が長い）は中央から遠く、紫色（波長が短い）は中央に近い位置に明線ができる。なお、$m=0$ では波長によらず $\theta=0$ であるから、中央の明線は白色のままである\ans
\ascboxZ{確認ポイント}\par
(3) の $\Delta y=2.5\unit{mm}$ は、(5) のずれ $12.5\unit{mm}$ のちょうど $\dfrac{1}{5}$ になっている。これは $(n-1)t=5\lambda$ という設定と整合する。
\begin{simple}[補足]{誤答しやすい点}
\begin{itemize}[leftmargin=1.2\zw,itemsep=0pt,topsep=0pt,parsep=0pt]
\item ガラス中で波長が変わるのを、ずれの計算に直接使おうとする。光路長 $(n-1)t$ で考えると簡単である。
\item 中央の明線のずれの向きを逆にする。光路長が増える側へ、同じ光路差になる点が動く。
\item (7) で $m=\pm4$ を数えて 9 本とする。$\sin\theta=1$ は観測されない。
\end{itemize}
\end{simple}
\end{multicols*}

% ---------------------------- 解答 大問4 ----------------------------
\begin{ascolorbox4A}[75点・30分]{大問4　解答・解説}
\textgt{分野タグ}　\#力のつり合い　\#状態方程式　\#熱力学第一法則　\#定圧・定積変化　\#分子運動論　\#p-Vグラフ\par
\textgt{検算}　済（$W=Mgh_0+p_0Sh_0$ の内訳、$Q_\mathrm{net}=W_\mathrm{net}=nRT_0$、$T_\mathrm{C}=T_\mathrm{A}$ を記号計算で確認した）\par
\textgt{答}　(1) $p_1=p_0+\dfrac{Mg}{S}$　(2) $h_0=\dfrac{nRT_0}{p_0S+Mg}$　(3) $T_\mathrm{B}=2T_0$　(4) $Q_\mathrm{AB}=\dfrac{5}{2}nRT_0$　(5) (i) $Mgh_0$、(ii) $p_0Sh_0$　(6) 解説参照　(7) $p_\mathrm{C}=\dfrac{p_1}{2}$、$Q_\mathrm{BC}=\dfrac{3}{2}nRT_0$（放出）　(8) 図 12　(9) $Q_\mathrm{net}=nRT_0$（解説参照）
\end{ascolorbox4A}
\begin{multicols*}{2}
\ascboxZ{着眼点}\par
ピストンが自由に動くとき、気体の圧力はピストンの力のつり合いで決まり、ゆっくり動かす間は一定（定圧変化）である。気体がした仕事は、ピストンを持ち上げる仕事と、大気を押しのける仕事に分けられる。圧力が変わらない理由は、分子運動論で説明できる。
\begin{ptbs}{KEY}[使う公式]
\begin{itemize}[leftmargin=1.2\zw,itemsep=0pt,topsep=0pt,parsep=0pt]
\item 状態方程式 $pV=nRT$、$U=\dfrac{3}{2}nRT$
\item 定圧変化 $W=p\Delta V=nR\Delta T$
\item 第一法則 $Q=\Delta U+W$
\item 圧力 $p=\dfrac{1}{3}\dfrac{N}{V}m\overline{v^2}$
\end{itemize}
\end{ptbs}
\ascboxZ{解答・解説}\par
\textbf{(1)}\ ピストンにはたらく力は、下向きに重力 $Mg$ と外気の力 $p_0S$、上向きに気体の力 $p_1S$ である。つり合いより
\[ p_1S=p_0S+Mg\ \ \therefore\ p_1=p_0+\frac{Mg}{S}\ansm \]

\textbf{(2)}\ 状態 A の状態方程式 $p_1\cdot Sh_0=nRT_0$ より
\[ h_0=\frac{nRT_0}{p_1S}=\frac{nRT_0}{p_0S+Mg}\ansm \]

\textbf{(3)}\ ゆっくり動くので圧力は常に $p_1$ で一定である。体積は絶対温度に比例し、$h_0\to2h_0$ で体積が $2$ 倍になるから
\[ T_\mathrm{B}=2T_0\ansm \]

\textbf{(4)}\ 内部エネルギーの増加は
\[ \Delta U=\frac{3}{2}nR(T_\mathrm{B}-T_0)=\frac{3}{2}nRT_0 \]
気体がした仕事は、(2) の状態方程式 $p_1Sh_0=nRT_0$ を用いて
\[ W=p_1\Delta V=p_1Sh_0=nRT_0 \]
第一法則より
\[ Q_\mathrm{AB}=\Delta U+W=\frac{5}{2}nRT_0\ansm \]
吸収した熱の $\dfrac{3}{5}$ が内部エネルギーに、$\dfrac{2}{5}$ が仕事に使われる。

\textbf{(5)}\ 気体がした仕事は
\[ W=p_1Sh_0=(p_0S+Mg)h_0=Mgh_0+p_0Sh_0 \]
である。(i) ピストンの重力による位置エネルギーの増加分は $Mgh_0$\ans\ (ii) 大気に対してした仕事は $p_0Sh_0$\ans\ 
ゆっくり動くので、ピストンの運動エネルギーの変化は $0$ である。

\textbf{(6)}\ A→B で温度が $2$ 倍になると、分子の平均の運動エネルギーが $2$ 倍になり、二乗平均速度は $\sqrt{2}$ 倍になる。一方、体積が $2$ 倍になるので、単位体積あたりの分子数は $\dfrac{1}{2}$ 倍になる。壁に衝突するとき 1 回あたりに壁へ与える力積は、分子の速さに比例して $\sqrt{2}$ 倍になる。壁の単位面積に単位時間あたり衝突する分子の数は、（分子数密度）$\times$（速さ）に比例し、$\dfrac{1}{2}\times\sqrt{2}=\dfrac{1}{\sqrt{2}}$ 倍になる。圧力はこの 2 つの積に比例するから $\sqrt{2}\times\dfrac{1}{\sqrt{2}}=1$ 倍となり、圧力は変わらない\ans
（$p=\dfrac{1}{3}\dfrac{N}{V}m\overline{v^2}$ で、$\dfrac{N}{V}$ が $\dfrac{1}{2}$ 倍、$\overline{v^2}$ が $2$ 倍と見ても同じ結論である。）

\textbf{(7)}\ B→C は体積一定であるから $\dfrac{p}{T}$ が一定である。
\[ p_\mathrm{C}=p_1\times\frac{T_0}{2T_0}=\frac{p_1}{2}\ansm \]
定積変化で $W=0$ であるから
\[ Q_\mathrm{BC}=\Delta U=\frac{3}{2}nR(T_0-2T_0)=-\frac{3}{2}nRT_0 \]
したがって、放出した熱量は $\dfrac{3}{2}nRT_0$ である\ans

\textbf{(8)}\ 体積を $V_0=Sh_0$ とおくと、A $(V_0,\,p_1)$、B $(2V_0,\,p_1)$、C $(2V_0,\,\dfrac{p_1}{2})$ である。A→B は横軸に平行な右向きの線分、B→C は縦軸に平行な下向きの線分である（図 12）\ans
\begin{center}
\begin{tikzpicture}[>={Stealth[length=1.8mm]},line width=0.6pt,x=1.35cm,y=1.1cm]
  \draw[->] (0,0) -- (3.0,0) node[right] {\small $V$};
  \draw[->] (0,0) -- (0,2.5) node[above] {\small $p$};
  \node[below left] at (0,0) {\small O};
  \draw[densely dashed] (1,0) node[below] {\small $V_0$} -- (1,2);
  \draw[densely dashed] (2,0) node[below] {\small $2V_0$} -- (2,2);
  \draw[densely dashed] (0,2) node[left] {\small $p_1$} -- (1,2);
  \draw[densely dashed] (0,1) node[left] {\small $\tfrac{p_1}{2}$} -- (2,1);
  \draw[densely dotted,domain=0.85:2.3,samples=40] plot ({\x},{2/\x});
  \draw[thick,->] (1,2) -- (1.6,2);
  \draw[thick] (1.6,2) -- (2,2);
  \draw[thick,->] (2,2) -- (2,1.4);
  \draw[thick] (2,1.4) -- (2,1);
  \fill (1,2) circle[radius=0.05] node[above left] {\small A};
  \fill (2,2) circle[radius=0.05] node[above right] {\small B};
  \fill (2,1) circle[radius=0.05] node[below right] {\small C};
\end{tikzpicture}\\[1pt]{\small 図12　点線は A, C を通る等温変化}
\end{center}

\textbf{(9)}\ 
\[ Q_\mathrm{net}=Q_\mathrm{AB}+Q_\mathrm{BC} \]
\[ =\frac{5}{2}nRT_0-\frac{3}{2}nRT_0=nRT_0\ansm \]
状態 C の温度は $T_0$ で、状態 A と等しい。内部エネルギーは温度だけで決まるから、A→B→C 全体で $\Delta U_\mathrm{net}=0$ である。第一法則より $Q_\mathrm{net}=\Delta U_\mathrm{net}+W_\mathrm{net}=W_\mathrm{net}$ となり、$Q_\mathrm{net}=W_\mathrm{net}$ が成り立つ\ans
実際、$W_\mathrm{net}=W_\mathrm{AB}+W_\mathrm{BC}=nRT_0+0=nRT_0$ で一致する。
\ascboxZ{確認ポイント}\par
(4) の $\dfrac{5}{2}nRT_0$ は、単原子分子の定圧モル比熱 $C_p=\dfrac{5}{2}R$ を用いて $nC_p(T_\mathrm{B}-T_0)$ としても求められる。(7) の $p_\mathrm{C}$ は、状態方程式 $p_\mathrm{C}\cdot2V_0=nRT_0$ からも $\dfrac{p_1}{2}$ になる。
\begin{simple}[補足]{誤答しやすい点}
\begin{itemize}[leftmargin=1.2\zw,itemsep=0pt,topsep=0pt,parsep=0pt]
\item (1) で外気圧 $p_0$ だけを気体の圧力としてしまう。ピストンの重さも加える。
\item (5) の仕事を $p_0S\cdot h_0$ だけにする。ピストンを持ち上げる $Mgh_0$ も含まれる。
\item (6) で「温度が上がるから圧力も上がる」とだけ書く。体積の増加による分子数密度の減少にも触れる。
\item (7) の $Q_\mathrm{BC}$ の符号を正にする。冷却は放熱で、$Q<0$ である。
\end{itemize}
\end{simple}
\end{multicols*}
% ==================================================================
%  解答・解説編　セット3（追加）
% ==================================================================
\clearpage
\begin{simple}[セット3（追加）]{解答・解説編　凡例}
各大問の冒頭の枠に、分野タグ・検算の状況・答の一覧を示す。解説は「着眼点 → 使う公式 → 解答・解説 → 確認ポイント → 誤答しやすい点」の順に述べる。
\end{simple}

% ---------------------------- 解答 大問1 ----------------------------
\begin{ascolorbox4A}[75点・30分]{大問1　解答・解説}
\textgt{分野タグ}　\#コンデンサー　\#電荷保存　\#エネルギー保存　\#電池のする仕事　\#誘電体\par
\textgt{検算}　済（全設問を数値で再計算し、エネルギー収支 $W=\Delta U'+H$、および $H=\tfrac{1}{2}(C_1+C_2)(E-V')^2$ を確認した）\par
\textgt{答}　(1) $Q_1=3.6\times10^{-5}\unit{C}$、$U_1=2.2\times10^{-4}\unit{J}$　(2) $V'=4.0\unit{V}$、$Q_1'=1.2\times10^{-5}\unit{C}$、$Q_2'=2.4\times10^{-5}\unit{C}$　(3) $\Delta U=1.4\times10^{-4}\unit{J}$、導線の抵抗でのジュール熱（と電磁波）　(4) $\Delta Q=7.2\times10^{-5}\unit{C}$、$W=8.6\times10^{-4}\unit{J}$　(5) $\Delta U'=5.8\times10^{-4}\unit{J}$、$H=2.9\times10^{-4}\unit{J}$　(6) $V''=7.2\unit{V}$　(7) $2.6\times10^{-4}\unit{J}$ 減少
\end{ascolorbox4A}
\begin{multicols*}{2}
\ascboxZ{着眼点}\par
コンデンサーの問題は、「電池がつながっているか（電圧が一定）」と「孤立しているか（電気量が一定）」を、操作のたびに見分けることが出発点になる。電池が切り離されたあとは、つながった極板全体の電気量の和が保存される。エネルギーは、静電エネルギーだけでは保存せず、ジュール熱や仕事とあわせて収支を考える。
\begin{ptbs}{KEY}[使う公式]
\begin{itemize}[leftmargin=1.2\zw,itemsep=0pt,topsep=0pt,parsep=0pt]
\item $Q=CV$、$U=\dfrac{1}{2}CV^2=\dfrac{Q^2}{2C}$
\item 並列：電圧が共通、合成容量 $C_1+C_2$
\item 電池がした仕事 $W=E\times$（移動した電気量）
\item 比誘電率 $\varepsilon_\mathrm{r}$ の誘電体で満たす：$C\to\varepsilon_\mathrm{r}C$
\end{itemize}
\end{ptbs}
\ascboxZ{解答・解説}\par
\textbf{(1)}\ $C_1$ には電圧 $E$ がかかる。
\[ Q_1=C_1E=3.0\times10^{-6}\times12=3.6\times10^{-5}\unit{C}\ansm \]
\[ U_1=\frac{1}{2}C_1E^2=\frac{1}{2}\times3.0\times10^{-6}\times12^2=2.16\times10^{-4}\unit{J} \]
有効数字 2 桁で $2.2\times10^{-4}\unit{J}$\ans

\textbf{(2)}\ $\mathrm{S}_1$ を開くと $C_1$ は孤立し、電気量 $Q_1$ が保存される。$\mathrm{S}_2$ を閉じると $C_1$、$C_2$ は並列になり、同じ電圧 $V'$ をもつ。電気量の和が保存されるから
\[ (C_1+C_2)V'=Q_1 \]
\[ V'=\frac{3.6\times10^{-5}}{9.0\times10^{-6}}=4.0\unit{V}\ansm \]
\[ Q_1'=C_1V'=1.2\times10^{-5}\unit{C} \]
\[ Q_2'=C_2V'=2.4\times10^{-5}\unit{C}\ansm \]
和は $3.6\times10^{-5}\unit{C}$ で、$Q_1$ に等しい。

\textbf{(3)}\ (2) の状態の静電エネルギーは
\[ U_2=\frac{1}{2}(C_1+C_2)V'^2 \]
\[ =\frac{1}{2}\times9.0\times10^{-6}\times4.0^2=7.2\times10^{-5}\unit{J} \]
であるから
\[ \Delta U=U_1-U_2=1.44\times10^{-4}\unit{J}\fallingdotseq1.4\times10^{-4}\unit{J}\ansm \]
失われたエネルギーは、電荷が導線を移動するときの抵抗によるジュール熱（と一部は電磁波）に変わった\ans
この値は、導線の抵抗の大きさによらない。

\textbf{(4)}\ $\mathrm{S}_1$ を閉じると、並列の 2 つのコンデンサーに電圧 $E$ がかかる。最終的な電気量の和は
\[ (C_1+C_2)E=9.0\times10^{-6}\times12=1.08\times10^{-4}\unit{C} \]
であり、操作の前は $3.6\times10^{-5}\unit{C}$ であったから、電池を通って移動した電気量は
\[ \Delta Q=1.08\times10^{-4}-3.6\times10^{-5}=7.2\times10^{-5}\unit{C}\ansm \]
\[ W=E\,\Delta Q=12\times7.2\times10^{-5} \]
\[ =8.64\times10^{-4}\unit{J}\fallingdotseq8.6\times10^{-4}\unit{J}\ansm \]

\textbf{(5)}\ 静電エネルギーの増加は
\[ \Delta U'=\frac{1}{2}(C_1+C_2)E^2-U_2 \]
\[ =6.48\times10^{-4}-7.2\times10^{-5}=5.76\times10^{-4}\unit{J} \]
有効数字 2 桁で $5.8\times10^{-4}\unit{J}$\ans\ 
電池がした仕事 $W$ が、静電エネルギーの増加とジュール熱になるから
\[ H=W-\Delta U' \]
\[ =8.64\times10^{-4}-5.76\times10^{-4}=2.88\times10^{-4}\unit{J}\ansm \]
有効数字 2 桁で $2.9\times10^{-4}\unit{J}$ である。

\textbf{(6)}\ $\mathrm{S}_1$ を開いたので、つながった極板の電気量の和 $1.08\times10^{-4}\unit{C}$ が保存される。$C_2$ は誘電体により $\varepsilon_\mathrm{r}C_2=12\,\mu\mathrm{F}$ になる。並列の合成容量は $3.0+12=15\,\mu\mathrm{F}$ であるから
\[ V''=\frac{1.08\times10^{-4}}{1.5\times10^{-5}}=7.2\unit{V}\ansm \]

\textbf{(7)}\ 操作の前後の静電エネルギーの合計は
\[ U_\text{前}=\frac{1}{2}\times9.0\times10^{-6}\times12^2=6.48\times10^{-4}\unit{J} \]
\[ U_\text{後}=\frac{1}{2}\times15\times10^{-6}\times7.2^2=3.888\times10^{-4}\unit{J} \]
\[ U_\text{後}-U_\text{前}=-2.592\times10^{-4}\unit{J} \]
有効数字 2 桁で、$2.6\times10^{-4}\unit{J}$ だけ減少する\ans\ 
誘電体は、静電気力によって極板の間に引き込まれる向きの力を受ける。誘電体を入れていく間に、この静電気力は誘電体に正の仕事をし、その分だけ静電エネルギーが減る。減った分は、誘電体の運動エネルギーとして現れるか、誘電体を静かに入れるときに外力（手）が受け止める。\ans
\ascboxZ{確認ポイント}\par
(5) のジュール熱は、$H=\dfrac{1}{2}(C_1+C_2)(E-V')^2=\dfrac{1}{2}\times9.0\times10^{-6}\times8.0^2=2.88\times10^{-4}\unit{J}$ としても求められ、先の値と一致する。
\begin{simple}[補足]{誤答しやすい点}
\begin{itemize}[leftmargin=1.2\zw,itemsep=0pt,topsep=0pt,parsep=0pt]
\item (2) で、電圧が $E$ のまま、あるいは電気量が $Q_1$ ずつ残ると考える。並列で電圧がそろい、電気量の和が保存される。
\item (4) の電池がした仕事を $\Delta U'$（静電エネルギーの増加）と等しいとする。半分はジュール熱になる。
\item (6) で、電池がつながっているときの考え方（$V''=E$）を使う。ここでは電池は切り離されている。
\end{itemize}
\end{simple}
\end{multicols*}

% ---------------------------- 解答 大問2 ----------------------------
\begin{ascolorbox4A}[75点・30分]{大問2　解答・解説}
\textgt{分野タグ}　\#運動量保存　\#動摩擦力　\#相対運動　\#熱への変換　\#v-tグラフ　\#弾性衝突\par
\textgt{検算}　済（時間・距離は記号計算し、$\ell_1$、$\ell_2$ はエネルギー損失と相対運動の 2 通りで一致を確認した。数値 $v_0=7.0\unit{m/s}$、$\mu'=0.50$ で $\ell_1=10/3\unit{m}$）\par
\textgt{答}　(1) $a=\mu'g$、$b=\dfrac{\mu'g}{2}$　(2) $t_1=\dfrac{2v_0}{3\mu'g}$、$v_1=\dfrac{v_0}{3}$　(3) $\ell_1=\dfrac{v_0^2}{3\mu'g}$、$3.3\unit{m}$　(4) 図 13　(5) A：$+\dfrac{v_0}{3}$、B：$-\dfrac{v_0}{3}$、$V=-\dfrac{v_0}{9}$（左向き）　(6) $\ell_2=\dfrac{4v_0^2}{27\mu'g}$　(7) $\ell_0\geqq\dfrac{13v_0^2}{27\mu'g}$
\end{ascolorbox4A}
\begin{multicols*}{2}
\ascboxZ{着眼点}\par
A と B の間にはたらく摩擦力は、作用・反作用の関係にある一対の力で、A と B の系にとっては内力である。したがって、床がなめらかなら、衝突のない間は A と B の運動量の和が保存される。「滑った距離」は、相対運動として、または失われた力学的エネルギー（摩擦熱）として求められる。壁との衝突では、外力（壁）がはたらくので、運動量は保存されない。
\begin{ptbs}{KEY}[使う公式]
\begin{itemize}[leftmargin=1.2\zw,itemsep=0pt,topsep=0pt,parsep=0pt]
\item 動摩擦力 $f'=\mu'N$、運動方程式 $ma=F$
\item 運動量保存 $m_1v_1+m_2v_2=\text{一定}$
\item 摩擦熱 $=\mu'N\times$（相対的に滑った距離）
\item 弾性衝突：壁に対する速さは変わらず、向きが逆になる
\end{itemize}
\end{ptbs}
\ascboxZ{解答・解説}\par
\textbf{(1)}\ A、B の間の動摩擦力の大きさは $\mu'mg$ である。A は左向きにこの力を受けて減速し、B は右向きにこの力を受けて加速する。
\[ a=\frac{\mu'mg}{m}=\mu'g\ansm \]
\[ b=\frac{\mu'mg}{2m}=\frac{\mu'g}{2}\ansm \]

\textbf{(2)}\ 時刻 $t$ での速度は、A が $v_0-\mu'gt$、B が $\dfrac{\mu'g}{2}t$ である。これらが等しくなる時刻が $t_1$ である。
\[ v_0-\mu'gt_1=\frac{\mu'g}{2}t_1\ \ \therefore\ t_1=\frac{2v_0}{3\mu'g}\ansm \]
\[ v_1=\frac{\mu'g}{2}t_1=\frac{v_0}{3}\ansm \]
運動量保存 $mv_0=(m+2m)v_1$ からも $v_1=\dfrac{v_0}{3}$ となり、一致する。

\textbf{(3)}\ 失われた力学的エネルギーは
\[ \Delta E_1=\frac{1}{2}mv_0^2-\frac{1}{2}(3m)\left(\frac{v_0}{3}\right)^2=\frac{1}{3}mv_0^2 \]
これが摩擦熱 $\mu'mg\,\ell_1$ に等しい。
\[ \mu'mg\,\ell_1=\frac{1}{3}mv_0^2\ \ \therefore\ \ell_1=\frac{v_0^2}{3\mu'g}\ansm \]
$v_0=7.0\unit{m/s}$、$\mu'=0.50$ のとき
\[ \ell_1=\frac{49}{3\times0.50\times9.8}=\frac{10}{3}\fallingdotseq3.3\unit{m}\ansm \]
別解として、B から見た A の相対運動は、初速 $v_0$、減速の大きさ $a+b=\dfrac{3}{2}\mu'g$ であるから $\ell_1=\dfrac{v_0^2}{2\left(\frac{3}{2}\mu'g\right)}=\dfrac{v_0^2}{3\mu'g}$ となる。

\textbf{(4)}\ $0\le t\le t_1$ では、A の速度は傾き $-\mu'g$ の直線で $v_0$ から $v_1=\dfrac{v_0}{3}$ まで減り、B の速度は傾き $+\dfrac{\mu'g}{2}$ の直線で $0$ から $v_1$ まで増える。$t_1\le t\le2t_1$ では、2 つとも $v_1$ で一定である（図 13）\ans
\begin{center}
\begin{tikzpicture}
\begin{axis}[width=5.9cm,height=3.4cm,axis lines=middle,
  xlabel={$t$},ylabel={速度},
  xlabel style={at={(ticklabel* cs:1.0)},anchor=west},
  ylabel style={at={(ticklabel* cs:1.0)},anchor=south,font=\scriptsize},
  xmin=0,xmax=2.5,ymin=0,ymax=3.7,
  xtick={1,2},xticklabels={$t_1$,$2t_1$},
  ytick={1,3},yticklabels={$v_1\!=\!\tfrac{v_0}{3}$,$v_0$},
  tick label style={font=\scriptsize},label style={font=\scriptsize}]
  \addplot[thick] coordinates {(0,3) (1,1) (2,1)};
  \addplot[thick,densely dashed] coordinates {(0,0) (1,1) (2,1)};
  \draw[densely dotted] (axis cs:1,0) -- (axis cs:1,1);
  \node[right] at (axis cs:0.02,2.7) {\scriptsize A};
  \node[below right] at (axis cs:0.45,0.2) {\scriptsize B};
\end{axis}
\end{tikzpicture}\\[1pt]{\small 図13　実線：A、破線：B}
\end{center}

\textbf{(5)}\ 弾性衝突では、板の速さは変わらず向きが逆になるから、B の速度は $-\dfrac{v_0}{3}$ となる\ans\ 
衝突の間に A にはたらく摩擦力の力積は無視できるので、A の速度は $+\dfrac{v_0}{3}$ のままである\ans\ 
この衝突の間は壁から外力を受けるので、A と B の系の運動量は保存されないが、衝突後は再び保存される。右向きを正として
\[ m\cdot\frac{v_0}{3}+2m\cdot\left(-\frac{v_0}{3}\right)=(m+2m)V \]
\[ V=-\frac{v_0}{9}\ansm \]
すなわち、左向きに大きさ $\dfrac{v_0}{9}$ で、A と B は一体となって動く。

\textbf{(6)}\ 衝突の直後から速度が等しくなるまでに失われた力学的エネルギーは
\[ \Delta E_2=\frac{1}{2}(3m)\left(\frac{v_0}{3}\right)^2-\frac{1}{2}(3m)\left(\frac{v_0}{9}\right)^2=\frac{4}{27}mv_0^2 \]
これが摩擦熱 $\mu'mg\,\ell_2$ に等しい。
\[ \ell_2=\frac{4v_0^2}{27\mu'g}\ansm \]
別解として、衝突直後の A の B に対する相対速度は $\dfrac{v_0}{3}-\left(-\dfrac{v_0}{3}\right)=\dfrac{2v_0}{3}$ で、減速の大きさ $a+b=\dfrac{3}{2}\mu'g$ であるから $\ell_2=\dfrac{(2v_0/3)^2}{2\cdot\frac{3}{2}\mu'g}=\dfrac{4v_0^2}{27\mu'g}$ となり、一致する。

\textbf{(7)}\ A は板の左端から出発し、B に対して右向きに、1 回目に $\ell_1$、衝突後にさらに $\ell_2$ だけ滑る（A の B に対する相対速度は、どちらの段階でも右向きである）。その後は一体となって左へ動くので、これ以上は滑らない。したがって、板の右端から落ちない条件は
\[ \ell_0\geqq\ell_1+\ell_2=\frac{9v_0^2+4v_0^2}{27\mu'g}=\frac{13v_0^2}{27\mu'g}\ansm \]
\ascboxZ{確認ポイント}\par
全体の損失は $\Delta E_1+\Delta E_2=\left(\dfrac{1}{3}+\dfrac{4}{27}\right)mv_0^2=\dfrac{13}{27}mv_0^2$ で、$\mu'mg(\ell_1+\ell_2)$ に等しい。壁は仕事をしない（弾性衝突）ので、失われるエネルギーは摩擦熱だけである。
\begin{simple}[補足]{誤答しやすい点}
\begin{itemize}[leftmargin=1.2\zw,itemsep=0pt,topsep=0pt,parsep=0pt]
\item 壁との衝突で、A と B の運動量が保存されるとして $V$ を求める。壁から外力を受ける。
\item 衝突後の B の速度の符号を誤る。右向きを正とすると、$-\dfrac{v_0}{3}$ である。
\item 「滑った距離」を、床に対する A の移動距離とする。板に対する相対的な距離で、摩擦熱とつながる。
\item (7) で、$\ell_1$ だけを答える。壁との衝突後の滑りも加える。
\end{itemize}
\end{simple}
\end{multicols*}

% ---------------------------- 解答 大問3 ----------------------------
\begin{ascolorbox4A}[75点・30分]{大問3　解答・解説}
\textgt{分野タグ}　\#電場と磁場　\#ローレンツ力　\#速度選択器　\#荷電粒子の円運動　\#質量分析\par
\textgt{検算}　済（数値は $v=1.0\times10^{4}\unit{m/s}$、$r=1.0\unit{cm}$、$\mathrm{PP}'=2.5\unit{mm}$ を再計算し、力の向きは外積で確認した）\par
\textgt{答}　(1) 奥向き、$v=\dfrac{V}{dB_1}$　(2) $1.0\times10^{4}\unit{m/s}$　(3) $r=\dfrac{mv}{qB_2}$、$m=\dfrac{qB_1B_2d\cdot\mathrm{OP}}{2V}$　(4) $1.0\times10^{-2}\unit{m}$　(5) $2.5\times10^{-3}\unit{m}$　(6) 解説参照　(7) $3.1\times10^{-6}\unit{s}$　(8) $5.0\unit{V}$
\end{ascolorbox4A}
\begin{multicols*}{2}
\ascboxZ{着眼点}\par
速度選択器では、電場による力と磁場による力（ローレンツ力）がつり合う粒子だけが直進する。このつり合いの式には質量が現れないので、質量や電気量によらず、速さ $v=\dfrac{E}{B_1}$ の粒子だけが通過する。そのあとは、一様な磁場の中で、ローレンツ力が向心力となる等速円運動である。
\begin{ptbs}{KEY}[使う公式]
\begin{itemize}[leftmargin=1.2\zw,itemsep=0pt,topsep=0pt,parsep=0pt]
\item 電場 $E=\dfrac{V}{d}$、電気力 $F=qE$
\item ローレンツ力 $f=qvB$（$\vec{v}\perp\vec{B}$）
\item 円運動 $m\dfrac{v^2}{r}=qvB$、周期 $T=\dfrac{2\pi m}{qB}$
\end{itemize}
\end{ptbs}
\ascboxZ{解答・解説}\par
\textbf{(1)}\ 上の極板が高電位であるから、電場は下向きで、正の電荷の粒子は下向きの電気力 $qE$ を受ける。直進するには、ローレンツ力が上向きでなければならない。粒子は右向きに動くので、$\vec{v}\times\vec{B}$ が上向きになる $\vec{B}$ の向きを考えると、磁場 $B_1$ は紙面に垂直で奥向きである\ans
つり合いの式は、$E=\dfrac{V}{d}$ を用いて
\[ qE=qvB_1\ \ \therefore\ v=\frac{E}{B_1}=\frac{V}{dB_1}\ansm \]

\textbf{(2)}\ 
\[ E=\frac{40}{2.0\times10^{-2}}=2.0\times10^{3}\unit{V/m} \]
\[ v=\frac{2.0\times10^{3}}{0.20}=1.0\times10^{4}\unit{m/s}\ansm \]

\textbf{(3)}\ ローレンツ力が向心力になる。
\[ m\frac{v^2}{r}=qvB_2\ \ \therefore\ r=\frac{mv}{qB_2}\ansm \]
$\mathrm{OP}=2r$ であるから $r=\dfrac{\mathrm{OP}}{2}$ で、これと $v=\dfrac{V}{dB_1}$ を $m=\dfrac{qB_2r}{v}$ に代入すると
\[ m=\frac{qB_2}{v}\cdot\frac{\mathrm{OP}}{2}=\frac{qB_1B_2d\cdot\mathrm{OP}}{2V}\ansm \]

\textbf{(4)}\ 
\[ r=\frac{1.6\times10^{-26}\times1.0\times10^{4}}{1.6\times10^{-19}\times0.10}=1.0\times10^{-2}\unit{m}\ansm \]

\textbf{(5)}\ 速度選択器を通過した粒子の速さは、質量によらず $v$ である。$r\propto m$ であるから
\[ r'=\frac{1.8}{1.6}\times1.0\times10^{-2}=1.125\times10^{-2}\unit{m} \]
\[ \mathrm{PP}'=2(r'-r) \]
\[ =2\times0.125\times10^{-2}=2.5\times10^{-3}\unit{m}\ansm \]

\textbf{(6)}\ 速さ $v'>v$ の粒子には、電気力 $qE$ は変わらないが、ローレンツ力 $qv'B_1$ が $qE$ より大きくなる。このため合力はローレンツ力の向き（上向き）になり、粒子は直進せずに上へ曲がる。したがってスリット $\mathrm{S}_2$ の位置に達することができず、極板や壁に当たる\ans

\textbf{(7)}\ 半周に要する時間は周期の半分である。
\[ t=\frac{1}{2}\cdot\frac{2\pi m}{qB_2}=\frac{\pi m}{qB_2}=\frac{\pi\times1.6\times10^{-26}}{1.6\times10^{-19}\times0.10}=\pi\times10^{-6}\unit{s} \]
有効数字 2 桁で $3.1\times10^{-6}\unit{s}$\ans\ 
この時間は、粒子の速さによらない。

\textbf{(8)}\ 電場が粒子にする仕事 $qV_0$ が運動エネルギーになる。
\[ qV_0=\frac{1}{2}mv^2 \]
\[ V_0=\frac{mv^2}{2q} \]
\[ =\frac{1.6\times10^{-26}\times(1.0\times10^{4})^2}{2\times1.6\times10^{-19}}=5.0\unit{V}\ansm \]
\ascboxZ{確認ポイント}\par
(5) では、粒子が描く半円は点 $\mathrm{P}$ より外側（下方）に延びる。質量が大きい粒子ほど半径が大きいので、$\mathrm{P}'$ は $\mathrm{P}$ より遠い。(3) の $m$ の式に $\mathrm{OP}=2r=2.0\times10^{-2}\unit{m}$ を代入すると、$m=1.6\times10^{-26}\unit{kg}$ が得られ、(4) と整合する。
\begin{simple}[補足]{誤答しやすい点}
\begin{itemize}[leftmargin=1.2\zw,itemsep=0pt,topsep=0pt,parsep=0pt]
\item (1) で、電場の向きと磁場の向きを同じ向きにそろえて考える。電気力とローレンツ力は逆向きでなければならない。
\item (5) で、質量が違えば速さも違うと考える。速度選択器を通る速さは質量によらない。
\item (3) の $\mathrm{OP}$ を $r$ とする。$\mathrm{OP}$ は半円の直径で、$2r$ である。
\end{itemize}
\end{simple}
\end{multicols*}

% ---------------------------- 解答 大問4 ----------------------------
\begin{ascolorbox4A}[75点・30分]{大問4　解答・解説}
\textgt{分野タグ}　\#単振動　\#ばね振り子　\#垂直抗力　\#物体の分離　\#鉛直投げ上げ\par
\textgt{検算}　済（$x_0=0.20\unit{m}$、$T=2\pi/7\unit{s}$、$v_1=1.05\unit{m/s}$ を、単振動のエネルギー保存と、放した位置から自然の長さまでの力学的エネルギー保存の 2 通りで確認した）\par
\textgt{答}　(1) $0.20\unit{m}$　(2) 解説参照、$T=\dfrac{2\pi}{7}\unit{s}\fallingdotseq0.90\unit{s}$　(3) $33\unit{N}$　(4) ばねが自然の長さになったとき、$v_1=1.05\unit{m/s}$　(5) $5.6\times10^{-2}\unit{m}$　(6) $d\leqq\dfrac{(M+m)g}{k}=0.20\unit{m}$　(7) $N=14.7-73.5x\ [\mathrm{N}]$（図 14）
\end{ascolorbox4A}
\begin{multicols*}{2}
\ascboxZ{着眼点}\par
重ねた物体が離れるかどうかは、上の物体が板から受ける垂直抗力 $N$ が $0$ になるかどうかで決まる。板と B が一体で単振動をしている間は、B の運動方程式から $N$ を位置の関数として表せる。$N=0$ となる位置が「離れる位置」で、そこから先は B だけが重力を受けて動く。
\begin{ptbs}{KEY}[使う公式]
\begin{itemize}[leftmargin=1.2\zw,itemsep=0pt,topsep=0pt,parsep=0pt]
\item 単振動 $a=-\omega^2x$、$\omega=\sqrt{\dfrac{k}{m_\text{全}}}$、$T=\dfrac{2\pi}{\omega}$
\item 単振動のエネルギー $\dfrac{1}{2}m_\text{全}v^2+\dfrac{1}{2}kx^2=\dfrac{1}{2}kd^2$
\item 離れる条件 $N=0$、投げ上げ $h=\dfrac{v^2}{2g}$
\end{itemize}
\end{ptbs}
\ascboxZ{解答・解説}\par
\textbf{(1)}\ つり合いでは、ばねの弾性力が板と B の重力とつり合う。
\[ kx_0=(M+m)g \]
\[ x_0=\frac{(M+m)g}{k}=\frac{2.0\times9.8}{98}=0.20\unit{m}\ansm \]

\textbf{(2)}\ つり合いの位置から上向きに $x$ だけ離れた位置では、ばねの縮みは $x_0-x$ であるから、板と B にはたらく合力は
\[ F=k(x_0-x)-(M+m)g=-kx \]
となる（$kx_0=(M+m)g$ を用いた）。合力が変位に比例し、常に変位と逆向きなので、板と B の運動は単振動である\ans\ 
\[ \omega=\sqrt{\frac{k}{M+m}}=\sqrt{\frac{98}{2.0}}=7.0\unit{rad/s} \]
\[ T=\frac{2\pi}{\omega}=\frac{2\pi}{7.0}\fallingdotseq0.90\unit{s}\ansm \]

\textbf{(3)}\ 放した直後は $x=-d=-0.25\unit{m}$ で、加速度は上向きに
\[ \omega^2d=49\times0.25=12.25\unit{m/s^2} \]
である。B の運動方程式（上向きを正）は $N_0-mg=m\omega^2d$ であるから
\[ N_0=m(g+\omega^2d)=1.5\times(9.8+12.25)=33.075\unit{N} \]
有効数字 2 桁で $33\unit{N}$\ans

\textbf{(4)}\ 位置 $x$ での加速度は $-\omega^2x$ であるから、B の運動方程式 $N-mg=m(-\omega^2x)$ より
\[ N=m(g-\omega^2x) \]
$N=0$ となる位置は
\[ x=\frac{g}{\omega^2}=\frac{(M+m)g}{k}=x_0=0.20\unit{m} \]
この位置では、ばねの縮み $x_0-x$ が $0$ である。したがって、B が離れるのは、ばねが自然の長さになったとき\ans\ 
その瞬間の速さは、単振動のエネルギー保存（振幅 $d$）から
\[ \frac{1}{2}(M+m)v_1^2+\frac{1}{2}kx_0^2=\frac{1}{2}kd^2 \]
\[ v_1^2=\omega^2(d^2-x_0^2)=49\times(0.25^2-0.20^2)=1.1025 \]
\[ v_1=1.05\unit{m/s}\ansm \]
別解として、放した位置（縮み $0.45\unit{m}$）から自然の長さまでの力学的エネルギー保存 $\dfrac{1}{2}k(0.45)^2=\dfrac{1}{2}(M+m)v_1^2+(M+m)g(0.45)$ からも同じ値が得られる。

\textbf{(5)}\ 離れたあとの B は、速さ $v_1$ で鉛直に投げ上げられた物体と同じ運動をする。
\[ h=\frac{v_1^2}{2g}=\frac{1.1025}{2\times9.8}\fallingdotseq5.6\times10^{-2}\unit{m}\ansm \]

\textbf{(6)}\ 離れないためには、振動の最高点 $x=d$ でも $N\geqq0$ であればよい。
\[ N(d)=m(g-\omega^2d)\geqq0 \]
\[ \therefore\ d\leqq\frac{g}{\omega^2}=\frac{(M+m)g}{k}\ansm \]
数値では $\dfrac{2.0\times9.8}{98}=0.20\unit{m}$ である。つまり $d\leqq x_0$ ということで、(4) の「離れる位置」が振動の範囲の外にあれば離れない。

\textbf{(7)}\ (4) より
\[ N=m(g-\omega^2x)=mg-\frac{mk}{M+m}x \]
\[ =14.7-73.5\,x\ \ [\mathrm{N}]\ansm \]
範囲は $-0.25\le x\le0.20$（$x=0.20$ で B が離れる）である。両端の値は、$x=-0.25\unit{m}$ で $N=14.7+18.4\fallingdotseq33\unit{N}$、$x=0.20\unit{m}$ で $N=0$ である。$x=0$ では $N=mg=14.7\unit{N}$ となる。グラフは、右下がりの直線である（図 14）\ans
\begin{center}
\begin{tikzpicture}
\begin{axis}[width=5.9cm,height=3.5cm,axis lines=middle,
  xlabel={$x\,[\mathrm{m}]$},ylabel={$N\,[\mathrm{N}]$},
  xlabel style={at={(ticklabel* cs:1.0)},anchor=west,font=\scriptsize},
  ylabel style={at={(ticklabel* cs:1.0)},anchor=south,font=\scriptsize},
  xmin=-0.34,xmax=0.32,ymin=0,ymax=42,
  xtick={-0.25,0.20},xticklabels={$-0.25$,$0.20$},
  ytick={14.7,33},yticklabels={$14.7$,$33$},
  tick label style={font=\scriptsize},label style={font=\scriptsize}]
  \addplot[thick] coordinates {(-0.25,33.075) (0.20,0)};
  \draw[densely dotted] (axis cs:-0.25,0) -- (axis cs:-0.25,33.075) -- (axis cs:0,33.075);
  \draw[densely dotted] (axis cs:0,14.7) -- (axis cs:0,14.7);
  \fill (axis cs:-0.25,33.075) circle[radius=1.2pt];
  \fill (axis cs:0.20,0) circle[radius=1.2pt];
\end{axis}
\end{tikzpicture}\\[1pt]{\small 図14}
\end{center}
\ascboxZ{確認ポイント}\par
$x=0.20\unit{m}$ で $N=14.7-73.5\times0.20=0$ となり、(4) の離れる位置と一致する。$d=0.25\unit{m}>x_0=0.20\unit{m}$ であるから、(6) の条件を満たさず、B は板から離れる。
\begin{simple}[補足]{誤答しやすい点}
\begin{itemize}[leftmargin=1.2\zw,itemsep=0pt,topsep=0pt,parsep=0pt]
\item 板だけの質量 $M$ で $\omega$ を求める。一体で運動している間は、質量は $M+m$ である。
\item B が離れる位置を、振動の最高点（$x=d$）とする。離れるのは $N=0$ となる位置で、ばねが自然の長さのときである。
\item (3) で $N_0=mg$ とする。加速度があるので $N_0=m(g+\omega^2d)$ である。
\item (4) の速さを、$\dfrac{1}{2}kd^2=\dfrac{1}{2}(M+m)v_1^2$ として求める。重力の位置エネルギーも考慮する（単振動のエネルギー保存の式を使えば自動的に含まれる）。
\end{itemize}
\end{simple}
\end{multicols*}
\end{document}
